% Generated mechanically from the frozen moduli.tex target.
% Do not edit this generated file; edit the bound locale source instead.

% OK, start here.
%

\setcounter{chapter}{107}
\chapter{模叠}


\phantomsection
\label{moduli-section-phantom}






\section{引言}
\label{moduli-section-introduction}

\noindent
本章验证若干模空间与模叠的基本性质，其中包括
$\mathit{Hom}$、$\mathit{Isom}$、$\Cohstack_{X/B}$、
$\Quotfunctor_{\mathcal{F}/X/B}$、$\Hilbfunctor_{X/B}$、
$\Picardstack_{X/B}$、$\Picardfunctor_{X/B}$、
$\mathit{Mor}_B(Z, X)$、$\Spacesstack'_{fp, flat, proper}$、
$\Polarizedstack$ 和 $\Complexesstack_{X/B}$。
% P12-ERR-0190: see ../control/ERRATA_CANDIDATES.jsonl
在适当假设下，我们已经证明这些对象是代数空间或代数叠；见
《Quot》中的章节
\ref{quot-section-hom}、
\ref{quot-section-isom}、
\ref{quot-section-stack-coherent-sheaves}、
\ref{quot-section-not-flat}、
\ref{quot-section-quot}、
\ref{quot-section-hilb}、
\ref{quot-section-picard-stack}、
\ref{quot-section-picard-functor}、
\ref{quot-section-relative-morphisms}、
\ref{quot-section-stack-of-spaces}、
\ref{quot-section-polarized} 和
\ref{quot-section-moduli-complexes}。
曲线叠记为 $\textit{Curves}$，它在《Quot》的章节
\ref{quot-section-curves} 中引入，并将在曲线模空间一章中讨论；见
《曲线模空间》中的章节
\ref{moduli-curves-section-stack-curves}。

\medskip\noindent
在某种意义上，本章沿着 Grothendieck 讲义
\cite{Gr-I}、
\cite{Gr-II}、
\cite{Gr-III}、
\cite{Gr-IV}、
\cite{Gr-V} 和
\cite{Gr-VI} 的思路展开。







\section{约定与语言滥用}
\label{moduli-section-conventions}

\noindent
我们继续采用《叠的性质》中的章节
\ref{stacks-properties-section-conventions} 所引入的约定和语言滥用。
除非另有说明，基概形均为 $\Spec(\mathbf{Z})$。






\section{Hom 与 Isom 的性质}
\label{moduli-section-hom-isom}

\noindent
设 $f : X \to B$ 是有限表示的代数空间态射。假设
$\mathcal{F}$ 和 $\mathcal{G}$ 是拟凝聚
$\mathcal{O}_X$-模。若 $\mathcal{G}$ 有限表示，在 $B$ 上平坦，
且其支撑在 $B$ 上真，则由
$$
T/B \longmapsto \Hom_{\mathcal{O}_{X_T}}(\mathcal{F}_T, \mathcal{G}_T)
$$
定义的函子 $\mathit{Hom}(\mathcal{F}, \mathcal{G})$
是 $B$ 上的仿射代数空间。若 $\mathcal{F}$ 有限表示，则
$\mathit{Hom}(\mathcal{F}, \mathcal{G}) \to B$
有限表示。见《Quot》中的命题 \ref{quot-proposition-hom}。

\medskip\noindent
若 $\mathcal{F}$ 和 $\mathcal{G}$ 都有限表示、在 $B$ 上平坦，
且其支撑在 $B$ 上真，则子函子
$$
\mathit{Isom}(\mathcal{F}, \mathcal{G}) \subset
\mathit{Hom}(\mathcal{F}, \mathcal{G})
$$
是 $B$ 上仿射且有限表示的代数空间。见《Quot》中的命题
\ref{quot-proposition-isom}。




\section{凝聚层的叠的性质}
\label{moduli-section-stack-coherent-sheaves}

\noindent
设 $f : X \to B$ 是分离且有限表示的代数空间态射。则参数化具有真支撑
的凝聚模平坦族的叠 $\Cohstack_{X/B}$ 是代数的。见《Quot》中的定理
\ref{quot-theorem-coherent-algebraic-general}。

\begin{lemma}
\label{moduli-lemma-coherent-diagonal-affine-fp}
$\Cohstack_{X/B}$ 在 $B$ 上的对角态射仿射且有限表示。
\end{lemma}

\begin{proof}
《Quot》中的引理 \ref{quot-lemma-coherent-diagonal}
已经证明该对角态射可由代数空间表示。从其证明可知，我们需要证明：
对一对有限表示的 $\mathcal{O}_{X_T}$-模
$\mathcal{F}$、$\mathcal{G}$，若它们在 $T$ 上平坦且支撑在 $T$
上真，则
$\mathit{Isom}(\mathcal{F}, \mathcal{G}) \to T$
仿射且有限表示。这已在章节 \ref{moduli-section-hom-isom} 中讨论。
\end{proof}


\begin{lemma}
\label{moduli-lemma-coherent-qs-lfp}
态射 $\Cohstack_{X/B} \to B$ 拟分离且局部有限表示。
\end{lemma}

\begin{proof}
要检验 $\Cohstack_{X/B} \to B$ 拟分离，需要证明其对角态射拟紧且
拟分离。这立即由引理 \ref{moduli-lemma-coherent-diagonal-affine-fp} 得出。
要证明 $\Cohstack_{X/B} \to B$ 局部有限表示，需要证明
$\Cohstack_{X/B} \to B$ 保持极限；见《叠的极限》中的命题
\ref{stacks-limits-proposition-characterize-locally-finite-presentation}。
这由《Quot》中的引理 \ref{quot-lemma-coherent-limits} 得出
（略去一个小细节）。
\end{proof}

\begin{lemma}
\label{moduli-lemma-coherent-existence-part}
再假设 $X \to B$ 为真且有限表示。则
$\Cohstack_{X/B} \to B$ 满足赋值判据的存在性部分
（《叠的态射》中的定义
\ref{stacks-morphisms-definition-existence}）。
\end{lemma}

\begin{proof}
作基变换后，问题立即约化为下述问题：给定分式域为 $K$ 的赋值环
$R$、一个在 $R$ 上真的代数空间 $X$，以及一个凝聚
$\mathcal{O}_{X_K}$-模 $\mathcal{F}_K$，证明存在一个有限表示的
$\mathcal{O}_X$-模 $\mathcal{F}$，它在 $R$ 上平坦，且其泛纤维为
$\mathcal{F}_K$。
注意，由《空间上的平坦性》中的定理
\ref{spaces-flat-theorem-finite-type-flat}，任意在 $R$ 上平坦的
有限型拟凝聚 $\mathcal{O}_X$-模 $\mathcal{F}$ 都有限表示。
% P12-ERR-0191: see ../control/ERRATA_CANDIDATES.jsonl
以 $j : X_K \to X$ 表示泛纤维的嵌入。由于它是仿射态射
$\Spec(K) \to \Spec(R)$ 的基变换，态射 $j$ 仿射。因此
$j_*\mathcal{F}_K$ 拟凝聚。把它写成
$$
j_*\mathcal{F}_K = \colim \mathcal{F}_i
$$
其中右边是它的有限型拟凝聚 $\mathcal{O}_X$-子模的滤过余极限；
见《空间的极限》中的引理
\ref{spaces-limits-lemma-directed-colimit-finite-type}。
由于 $j_*\mathcal{F}_K$ 是 $X$ 上的 $K$-向量空间层，它在
$\Spec(R)$ 上平坦。因此每个 $\mathcal{F}_i$ 都在 $R$ 上平坦；
这是因为在赋值环上平坦等价于无挠
（《代数进阶》中的引理
\ref{more-algebra-lemma-valuation-ring-torsion-free-flat}），
而子模继承无挠性。
最后，需要证明对某个 $i$，映射
$j^*\mathcal{F}_i \to \mathcal{F}_K$
是同构。由于
$j^*j_*\mathcal{F}_K = \mathcal{F}_K$（略去一个小细节），
且 $j^*$ 正合，可见对所有 $i$，映射
$j^*\mathcal{F}_i \to \mathcal{F}_K$ 都是单射。
由于 $j^*$ 与余极限交换，有
$\mathcal{F}_K = j^*j_*\mathcal{F}_K = \colim j^*\mathcal{F}_i$。
由于 $\mathcal{F}_K$ 凝聚（即有限表示），存在某个 $i$，使得
$j^*\mathcal{F}_i$ 包含 $X$ 的一个仿射 \'etale 覆盖上的全部
（有限多个）生成元。因此，当 $i$ 充分大时，
$j^*\mathcal{F}_i \to \mathcal{F}_K$ 满射。
\end{proof}
\begin{lemma}
\label{moduli-lemma-coherent-functorial}
设 $B$ 是代数空间。设 $\pi : X \to Y$ 是代数空间之间的拟有限态射，
其中两者均在 $B$ 上分离且有限表示。则 $\pi_*$ 诱导一个态射
$\Cohstack_{X/B} \to \Cohstack_{Y/B}$。
\end{lemma}

\begin{proof}
设 $(T \to B, \mathcal{F})$ 是 $\Cohstack_{X/B}$ 的一个对象。
我们断言：
\begin{enumerate}
\item[(a)] $(T \to B, \pi_{T, *}\mathcal{F})$ 是
$\Cohstack_{Y/B}$ 的对象；以及
\item[(b)] 对于 $T' \to T$，有
$\pi_{T', *}(X_{T'} \to X_T)^*\mathcal{F} =
(Y_{T'} \to Y_T)^*\pi_{T, *}\mathcal{F}$。
\end{enumerate}
(b) 保证这一构造按所需定义一个函子
$\Cohstack_{X/B} \to \Cohstack_{Y/B}$。

\medskip\noindent
设 $i : Z \to X_T$ 是由 $\mathcal{F}$ 的第零 Fitting 理想截出的
闭子空间
（《空间上的除子》中的章节
\ref{spaces-divisors-section-fitting-ideals}）。
% P12-ERR-0192: see ../control/ERRATA_CANDIDATES.jsonl
按假设，$Z \to B$ 为真（见《空间的导出范畴》中的章节
\ref{spaces-perfect-section-proper-over-base}）。
另一方面，$i$ 有限表示
（《空间上的除子》中的引理
\ref{spaces-divisors-lemma-fitting-ideal-of-finitely-presented}
以及《空间的态射》中的引理
\ref{spaces-morphisms-lemma-closed-immersion-finite-presentation}）。
存在有限型拟凝聚 $\mathcal{O}_Z$-模 $\mathcal{G}$，使得
$i_*\mathcal{G} = \mathcal{F}$
（《空间上的除子》中的引理
\ref{spaces-divisors-lemma-on-subscheme-cut-out-by-Fit-0}）。
事实上，由《空间上的下降》中的引理
\ref{spaces-descent-lemma-finite-finitely-presented-module}，
$\mathcal{G}$ 作为 $\mathcal{O}_Z$-模有限表示。
% P12-ERR-0193: see ../control/ERRATA_CANDIDATES.jsonl
注意 $\mathcal{G}$ 在 $B$ 上平坦；例如，这是因为
$\mathcal{G}$ 和 $\mathcal{F}$ 的茎相同
（《空间的态射》中的引理
\ref{spaces-morphisms-lemma-stalk-push-closed}）。
注意 $\pi_T \circ i : Z \to Y_T$ 作为拟有限态射的复合仍拟有限，并且
% P12-ERR-0194: see ../control/ERRATA_CANDIDATES.jsonl
$\pi_{T, *}\mathcal{F} = (\pi_T \circ i)_*\mathcal{G})$。
由于 $i$ 仿射，$i_*$ 的形成与基变换交换
（《空间的上同调》中的引理
\ref{spaces-cohomology-lemma-affine-base-change}）。
因此可把 $B$ 替换为 $T$，把 $X$ 替换为 $Z$，把
$\mathcal{F}$ 替换为 $\mathcal{G}$，并把 $Y$ 替换为 $Y_T$，
从而约化到下一段讨论的情形。

\medskip\noindent
假设 $X \to B$ 为真。则由《空间的态射》中的引理
\ref{spaces-morphisms-lemma-universally-closed-permanence}，
$\pi$ 为真，因而由《空间的态射进阶》中的引理
\ref{spaces-more-morphisms-lemma-characterize-finite}，它有限。
由于有限态射仿射，《空间的上同调》中的引理
\ref{spaces-cohomology-lemma-affine-base-change} 表明 (b) 成立。
另一方面，由《空间的态射》中的引理
\ref{spaces-morphisms-lemma-finite-presentation-permanence}，
$\pi$ 有限表示。因此，由《空间上的下降》中的引理
\ref{spaces-descent-lemma-finite-finitely-presented-module}，
$\pi_{T, *}\mathcal{F}$ 有限表示。
最后，例如利用《空间的上同调》中的引理
\ref{spaces-cohomology-lemma-stalk-push-finite} 考察茎，可知
$\pi_{T, *}\mathcal{F} $ 在 $B$ 上平坦。
\end{proof}

\begin{lemma}
\label{moduli-lemma-coherent-open}
设 $B$ 是代数空间。设 $\pi : X \to Y$ 是代数空间之间的开浸入，
其中两者均在 $B$ 上分离且有限表示。则引理
\ref{moduli-lemma-coherent-functorial} 中的态射
$\Cohstack_{X/B} \to \Cohstack_{Y/B}$ 是开浸入。
\end{lemma}

\begin{proof}
从略。提示：若 $\mathcal{F}$ 是 $\Cohstack_{Y/B}$ 在 $T$ 上的
一个对象，并且对 $t \in T$ 有
$\text{Supp}(\mathcal{F}_t) \subset |X_t|$，则对 $t$ 的某个邻域中的
$t' \in T$，同一结论也成立。
\end{proof}

\begin{lemma}
\label{moduli-lemma-coherent-closed}
设 $B$ 是代数空间。设 $\pi : X \to Y$ 是代数空间之间的闭浸入，
其中两者均在 $B$ 上分离且有限表示。则引理
\ref{moduli-lemma-coherent-functorial} 中的态射
$\Cohstack_{X/B} \to \Cohstack_{Y/B}$ 是闭浸入。
\end{lemma}

\begin{proof}
设 $\mathcal{I} \subset \mathcal{O}_Y$ 是把 $X$ 截为 $Y$ 的闭子空间的
理想层。回忆 $\pi_*$ 在拟凝聚 $\mathcal{O}_X$-模的范畴与被
$\mathcal{I}$ 零化的拟凝聚 $\mathcal{O}_Y$-模的范畴之间诱导等价；
见《空间的态射》中的引理
\ref{spaces-morphisms-lemma-i-star-equivalence}。
% P12-ERR-0195: see ../control/ERRATA_CANDIDATES.jsonl
对 $T \to B$ 作基变换并把 $\mathcal{I}$ 替换为理想层
$\mathcal{I}_T = \Im((Y_T \to Y)^*\mathcal{I} \to \mathcal{O}_{Y_T})$
后，作相应变更可得到同一结论。
分析引理 \ref{moduli-lemma-coherent-functorial} 的证明可知，态射
$\Cohstack_{X/B} \to \Cohstack_{Y/B}$ 的本质像恰好由对象
$\xi = (T \to B, \mathcal{F})$ 组成，其中
$\mathcal{F}$ 被 $\mathcal{I}_T$ 零化。
换言之，$\xi$ 位于本质像中，当且仅当乘法映射
$$
\mathcal{F} \otimes_{\mathcal{O}_{Y_T}} (Y_T \to Y)^*\mathcal{I}
\longrightarrow
\mathcal{F}
$$
为零，并且在任意进一步的基变换 $T' \to T$ 后也同样为零。
注意
$$
(Y_{T'} \to Y_T)^*(
\mathcal{F} \otimes_{\mathcal{O}_{Y_T}} (Y_T \to Y)^*\mathcal{I}) =
(Y_{T'} \to Y_T)^*\mathcal{F} \otimes_{\mathcal{O}_{Y_{T'}}}
(Y_{T'} \to Y)^*\mathcal{I})
$$
因此，由《空间上的平坦性》中的引理
\ref{spaces-flat-lemma-F-zero-closed-proper}，该乘法映射在 $T'$
上为零这一条件可由 $T$ 的某个闭子空间表示。
\end{proof}
\begin{situation}[数值不变量]
\label{moduli-situation-numerical}
设 $f : X \to B$ 如本节引言所述。设 $I$ 是一个集合，并对
$i \in I$ 设 $E_i \in D(\mathcal{O}_X)$ 为完美对象。给定
$\Cohstack_{X/B}$ 的一个对象 $(T \to B, \mathcal{F})$，
以 $E_{i, T}$ 表示 $E_i$ 到 $X_T$ 的导出拉回。$D(\mathcal{O}_T)$
中的对象
$$
K_i = Rf_{T, *}(E_{i, T} \otimes_{\mathcal{O}_{X_T}}^\mathbf{L} \mathcal{F})
$$
是完美的，且其形成与基变换交换；见《空间的导出范畴》中的引理
\ref{spaces-perfect-lemma-base-change-tensor-perfect}。
因此，由《空间的导出范畴》中的引理
\ref{spaces-perfect-lemma-chi-locally-constant}，函数
$$
\chi_i : |T| \longrightarrow \mathbf{Z},\quad
\chi_i(t) =
\chi(X_t, E_{i, t} \otimes_{\mathcal{O}_{X_t}}^\mathbf{L} \mathcal{F}_t) =
\chi(K_i \otimes_{\mathcal{O}_T}^\mathbf{L} \kappa(t))
$$
局部常值。设 $P : I \to \mathbf{Z}$ 是一个映射。考虑子叠
$$
\Cohstack^P_{X/B} \subset \Cohstack_{X/B}
$$
它由数值不变量与 $P$ 一致、具有真支撑的凝聚层平坦族组成。
更准确地，$\Cohstack_{X/B}$ 的一个对象
$(T \to B, \mathcal{F})$ 位于 $\Cohstack^P_{X/B}$ 中，当且仅当
对所有 $i \in I$ 和 $t \in T$ 都有 $\chi_i(t) = P(i)$。
\end{situation}

\begin{lemma}
\label{moduli-lemma-open-P}
在情形 \ref{moduli-situation-numerical} 中，叠
$\Cohstack^P_{X/B}$ 是代数的，而且
$$
\Cohstack^P_{X/B} \longrightarrow \Cohstack_{X/B}
$$
是平坦闭浸入。若 $I$ 有限，或 $B$ 局部 Noether，则
$\Cohstack^P_{X/B}$ 是 $\Cohstack_{X/B}$ 的既开又闭子叠。
\end{lemma}

\begin{proof}
若 $I$ 有限，这由函数 $t \mapsto \chi_i(t)$ 局部常值立即可知。
若 $I$ 无限，则写成
$$
I = \bigcup\nolimits_{I' \subset I\text{ 有限}} I'
$$
并记 $P' = P|_{I'}$。于是
$$
\Cohstack^P_{X/B} = \bigcap\nolimits_{I' \subset I\text{ 有限}}
\Cohstack^{P'}_{X/B}
$$
因此，$\Cohstack^P_{X/B}$ 总是代数叠，且态射
$\Cohstack^P_{X/B} \subset \Cohstack_{X/B}$ 总是平坦闭浸入，
但它可能不再是开子叠。（我们把例子的构造留给读者。）
不过，若 $B$ 局部 Noether，则由引理 \ref{moduli-lemma-coherent-qs-lfp}
和《叠的态射》中的引理
\ref{stacks-morphisms-lemma-locally-finite-type-locally-noetherian}，
$\Cohstack_{X/B}$ 也局部 Noether。因此，若
$U \to \Cohstack_{X/B}$ 是一个光滑满射，且 $U$ 是局部 Noether
概形，则既开又闭子叠 $\Cohstack^{P'}_{X/B}$ 的逆像在 $U$ 中的
交集是开集（因为局部 Noether 拓扑空间的连通分支是开的）。
这就证明了这种情形下的结论。
\end{proof}

\begin{lemma}
\label{moduli-lemma-finite-list-perfect-objects}
设 $f : X \to B$ 如本节引言所述。设
$E_1, \ldots, E_r \in D(\mathcal{O}_X)$ 为完美对象。
令 $I = \mathbf{Z}^{\oplus r}$，并考虑映射
$$
I \longrightarrow D(\mathcal{O}_X),\quad
(n_1, \ldots, n_r) \longmapsto
E_1^{\otimes n_1}
\otimes \ldots \otimes
E_r^{\otimes n_r}
$$
% P12-ERR-0196: see ../control/ERRATA_CANDIDATES.jsonl
设 $P : I \to \mathbf{Z}$ 是一个映射。则情形
\ref{moduli-situation-numerical} 中定义的
$\Cohstack^P_{X/B} \subset \Cohstack_{X/B}$
是既开又闭子叠。
\end{lemma}

\begin{proof}
可以在 $B$ 上 \'etale 局部地工作，故可设 $B$ 仿射。
在这种情形下，可以作绝对 Noether 约化；建议读者跳过此证明。
具体地，设 $B = \Spec(\Lambda)$。把
$\Lambda = \colim \Lambda_i$ 写成滤过余极限，其中每个
$\Lambda_i$ 都在 $\mathbf{Z}$ 上有限型。对某个 $i$，可以找到
一个分离且有限表示的代数空间态射
$X_i \to \Spec(\Lambda_i)$，其到 $\Lambda$ 的基变换为 $X$。
见《空间的极限》中的引理
\ref{spaces-limits-lemma-descend-finite-presentation} 和
\ref{spaces-limits-lemma-descend-separated-morphism}。
增大 $i$ 后，可进一步假设在 $D(\mathcal{O}_{X_i})$ 中存在完美对象
$E_{1, i}, \ldots, E_{r, i}$，它们到 $X$ 的导出拉回分别同构于
$E_1, \ldots, E_r$；见《空间的导出范畴》中的引理
\ref{spaces-perfect-lemma-perfect-on-limit}。
显然有 Cartesian 方块
$$
\xymatrix{
\Cohstack^P_{X/B} \ar[r] \ar[d] &
\Cohstack_{X/B} \ar[d] \\
\Cohstack^P_{X_i/\Spec(\Lambda_i)} \ar[r] &
\Cohstack_{X_i/\Spec(\Lambda_i)}
}
$$
因而可用引理 \ref{moduli-lemma-open-P} 完成证明。
\end{proof}

\begin{example}[具有固定 Hilbert 多项式的凝聚层]
\label{moduli-example-hilbert-polynomial}
设 $f : X \to B$ 如本节引言所述。设 $\mathcal{L}$ 是可逆
$\mathcal{O}_X$-模，并设
$P : \mathbf{Z} \to \mathbf{Z}$ 是数值多项式。于是可以考虑既开又闭的
代数子叠
$$
\Cohstack^P_{X/B} =
\Cohstack^{P, \mathcal{L}}_{X/B}
\subset \Cohstack_{X/B}
$$
它由数值不变量与 $P$ 一致、具有真支撑的凝聚层平坦族组成：
$\Cohstack_{X/B}$ 的一个对象
$(T \to B, \mathcal{F})$ 位于 $\Cohstack^P_{X/B}$ 中，当且仅当
$$
P(n) =
\chi(X_t, \mathcal{F}_t \otimes_{\mathcal{O}_{X_t}} \mathcal{L}_t^{\otimes n})
$$
对所有 $n \in \mathbf{Z}$ 和 $t \in T$ 成立。当然，这是情形
\ref{moduli-situation-numerical} 的特例，其中
% P12-ERR-0197: see ../control/ERRATA_CANDIDATES.jsonl
$I = \mathbf{Z} \to D(\mathcal{O}_X)$ 由
$n \mapsto \mathcal{L}^{\otimes n}$ 给出。
由引理 \ref{moduli-lemma-finite-list-perfect-objects}，这是既开又闭子叠。
由于函数
$n \mapsto
\chi(X_t, \mathcal{F}_t \otimes_{\mathcal{O}_{X_t}} \mathcal{L}_t^{\otimes n})$
总是数值多项式（《域上的空间》中的引理
\ref{spaces-over-fields-lemma-numerical-polynomial-from-euler}），
可得不交并分解
$$
\Cohstack_{X/B} = \coprod\nolimits_{P\text{ 为数值多项式}}
\Cohstack^P_{X/B}
$$
\end{example}






\section{Quot 的性质}
\label{moduli-section-quot}

\noindent
设 $f : X \to B$ 是分离且有限表示的代数空间态射，并设
$\mathcal{F}$ 是拟凝聚 $\mathcal{O}_X$-模。则
$\Quotfunctor_{\mathcal{F}/X/B}$ 是代数空间。若
$\mathcal{F}$ 有限表示，则
$\Quotfunctor_{\mathcal{F}/X/B} \to B$ 局部有限表示。
见《Quot》中的命题 \ref{quot-proposition-quot}。

\begin{lemma}
\label{moduli-lemma-quot-diagonal-closed}
$\Quotfunctor_{\mathcal{F}/X/B} \to B$ 的对角态射是闭浸入。
若 $\mathcal{F}$ 有限型，则该对角态射是有限表示的闭浸入。
\end{lemma}

\begin{proof}
设有一个概形 $T/B$，以及两个商
$\mathcal{F}_T \to \mathcal{Q}_i$（$i = 1, 2$），它们对应于
$\Quotfunctor_{\mathcal{F}/X/B}$ 在 $B$ 上的 $T$-值点。
以 $\mathcal{K}_1$ 表示第一个商的核，并以
$u : \mathcal{K}_1 \to \mathcal{Q}_2$ 表示复合。
由《空间上的平坦性》中的引理
\ref{spaces-flat-lemma-F-zero-closed-proper}，$T$ 中存在一个闭子空间，
使得 $T' \to T$ 通过它分解，当且仅当拉回 $u_{T'}$ 为零。
这就证明了对角态射是闭浸入。此外，若 $\mathcal{F}$ 有限型，则
$\mathcal{K}_1$ 有限型
（《位点上的模》中的引理
\ref{sites-modules-lemma-kernel-surjection-finite-onto-finite-presentation}），
并由同一引理可知该对角态射有限表示。
\end{proof}

\begin{lemma}
\label{moduli-lemma-quot-s-lfp}
态射 $\Quotfunctor_{\mathcal{F}/X/B} \to B$ 分离。
若 $\mathcal{F}$ 有限表示，则它还局部有限表示。
\end{lemma}

\begin{proof}
要检验 $\Quotfunctor_{\mathcal{F}/X/B} \to B$ 分离，需要证明其
对角态射为闭浸入。这由引理 \ref{moduli-lemma-quot-diagonal-closed} 得出。
第二个陈述是《Quot》中的命题 \ref{quot-proposition-quot} 的一部分。
\end{proof}

\begin{lemma}
\label{moduli-lemma-quot-existence-part}
假设 $X \to B$ 为真且有限表示，并且 $\mathcal{F}$ 是有限型拟凝聚模。
则 $\Quotfunctor_{\mathcal{F}/X/B} \to B$ 满足赋值判据的存在性部分
（《空间的态射》中的定义
\ref{spaces-morphisms-definition-valuative-criterion}）。
\end{lemma}

\begin{proof}
作基变换后，问题立即约化为下述问题：给定分式域为 $K$ 的赋值环
$R$、一个在 $R$ 上真的代数空间 $X$、一个有限型拟凝聚
$\mathcal{O}_X$-模 $\mathcal{F}$，以及一个凝聚商
$\mathcal{F}_K \to \mathcal{Q}_K$，证明存在一个商
$\mathcal{F} \to \mathcal{Q}$，其中 $\mathcal{Q}$ 是有限表示的
$\mathcal{O}_X$-模，在 $R$ 上平坦，且其泛纤维为 $\mathcal{Q}_K$。
注意，由《空间上的平坦性》中的定理
\ref{spaces-flat-theorem-finite-type-flat}，任意在 $R$ 上平坦的
有限型拟凝聚 $\mathcal{O}_X$-模 $\mathcal{F}$ 都有限表示。
先在仿射局部构造 $\mathcal{Q}$。

\medskip\noindent
仿射局部地，得到下述问题：设 $R \to A$ 是有限表示的环同态，
$M$ 是有限 $A$-模，且 $\varphi : M_K \to N_K$ 是一个 $A_K$-模商。
那么可以考虑
$$
L = \{x \in M \mid \varphi(x \otimes 1) = 0 \}
$$
% P12-ERR-0198: see ../control/ERRATA_CANDIDATES.jsonl
于是 $M \to M/L$ 是一个 $A$-模商，它作为 $R$-模无挠。因此，它作为
$R$-模平坦
（《代数进阶》中的引理
\ref{more-algebra-lemma-valuation-ring-torsion-free-flat}）。
由于 $M$ 作为 $A$-模有限，$L$ 也有限；由此得出
$L$ 作为 $A$-模有限表示（由上面的引用）。显然，$M/L$ 是满足
$(M/L)_K = N_K$ 的唯一这种商。

\medskip\noindent
上一段构造的唯一性保证这些商可以粘合，从而给出所需的
$\mathcal{Q}$。下面补充一些细节。选取一个满 \'etale 态射
$U \to X$，其中 $U$ 是仿射概形。用上述构造得到一个商
$\mathcal{F}|_U \to \mathcal{Q}_U$；它拟凝聚、在 $R$ 上平坦，
并在泛纤维上恢复 $\mathcal{Q}_K|U$。
由于 $X$ 分离，$U \times_X U$ 也是仿射概形，并且 \'etale 于
$X$。于是商
$\mathcal{F}|_{U \times_X U} \to \text{pr}_1^*\mathcal{Q}_U$ 和
$\mathcal{F}|_{U \times_X U} \to \text{pr}_2^*\mathcal{Q}_U$
由该构造的唯一性而一致。
% P12-ERR-0199: see ../control/ERRATA_CANDIDATES.jsonl
因此可把 $\mathcal{F}|_U \to \mathcal{Q}_U$ 下降为所需的满射
$\mathcal{F} \to \mathcal{Q}$
（《空间的性质》中的命题
\ref{spaces-properties-proposition-quasi-coherent}）。
\end{proof}

\begin{lemma}
\label{moduli-lemma-quot-functorial}
设 $B$ 是代数空间。设 $\pi : X \to Y$ 是代数空间之间的仿射拟有限态射，
且这些代数空间在 $B$ 上分离并有限表示。设 $\mathcal{F}$ 是拟凝聚
$\mathcal{O}_X$-模。则 $\pi_*$ 诱导态射
$\Quotfunctor_{\mathcal{F}/X/B} \to \Quotfunctor_{\pi_*\mathcal{F}/Y/B}$。
\end{lemma}

\begin{proof}
令 $\mathcal{G} = \pi_*\mathcal{F}$。由于 $\pi$ 仿射，对于任意
$B$ 上概形 $T$，由《空间的上同调》中的引理
\ref{spaces-cohomology-lemma-affine-base-change} 有
$\mathcal{G}_T = \pi_{T, *}\mathcal{F}_T$。此外，$\pi_T$ 仿射，
所以 $\pi_{T, *}$ 正合，并把商变为商。
% P12-ERR-0200: see ../control/ERRATA_CANDIDATES.jsonl
注意，拟凝聚商 $\mathcal{F}_T \to \mathcal{Q}$ 定义了
$\Quotfunctor_{X/B}$ 的一个点，当且仅当 $\mathcal{Q}$ 定义了
$\Cohstack_{X/B}$ 在 $T$ 上的一个对象（对 $\mathcal{G}$ 和 $Y$
也同样）。我们已在引理 \ref{moduli-lemma-coherent-functorial} 中看到，
$\pi_*$ 诱导态射 $\Cohstack_{X/B} \to \Cohstack_{Y/B}$。
因此，若 $\mathcal{F}_T \to \mathcal{Q}$ 属于
$\Quotfunctor_{\mathcal{F}/X/B}(T)$，则
$\mathcal{G}_T \to \pi_{T, *}\mathcal{Q}$ 属于
$\Quotfunctor_{\mathcal{G}/Y/B}(T)$。
\end{proof}

\begin{lemma}
\label{moduli-lemma-quot-open}
设 $B$ 是代数空间。设 $\pi : X \to Y$ 是代数空间之间的仿射开浸入，
且这些代数空间在 $B$ 上分离并有限表示。设 $\mathcal{F}$ 是拟凝聚
$\mathcal{O}_X$-模。则引理 \ref{moduli-lemma-quot-functorial} 的态射
$\Quotfunctor_{\mathcal{F}/X/B} \to
\Quotfunctor_{\pi_*\mathcal{F}/Y/B}$ 是开浸入。
\end{lemma}

\begin{proof}
略。提示：若 $(\pi_*\mathcal{F})_T \to \mathcal{Q}$ 是
$\Quotfunctor_{\pi_*\mathcal{F}/Y/B}(T)$ 的元素，并且对
$t \in T$ 有 $\text{Supp}(\mathcal{Q}_t) \subset |X_t|$，
那么对 $t$ 的某个邻域中的 $t' \in T$ 也有同一性质。
\end{proof}

\begin{lemma}
\label{moduli-lemma-quot-better-open}
设 $B$ 是代数空间。设 $j : X \to Y$ 是代数空间之间的开浸入，
且这些代数空间在 $B$ 上分离并有限表示。设 $\mathcal{G}$ 是拟凝聚
$\mathcal{O}_Y$-模，并令 $\mathcal{F} = j^*\mathcal{G}$。则存在
$B$ 上代数空间的开浸入
$$
\Quotfunctor_{\mathcal{F}/X/B}
\longrightarrow
\Quotfunctor_{\mathcal{G}/Y/B}
$$
\end{lemma}

\begin{proof}
若 $\mathcal{F}_T \to \mathcal{Q}$ 是
$\Quotfunctor_{\mathcal{F}/X/B}(T)$ 的元素，则可以考虑
$\mathcal{G}_T \to j_{T, *}\mathcal{F}_T \to j_{T, *}\mathcal{Q}$。
考察茎可知它是满射。由引理 \ref{moduli-lemma-coherent-functorial}，
% P12-ERR-0201: see ../control/ERRATA_CANDIDATES.jsonl
$j_{T, *}\mathcal{Q}$ 有限表示、在 $B$ 上平坦，且其支撑在 $B$ 上真。
因此得到 $\Quotfunctor_{\mathcal{G}/Y/B}$ 的一个 $T$-值点。
这就定义了本引理的态射。略去它是开浸入的证明。提示：若
$\mathcal{G}_T \to \mathcal{Q}$ 是
$\Quotfunctor_{\mathcal{G}/Y/B}(T)$ 的元素，并且对 $t \in T$ 有
$\text{Supp}(\mathcal{Q}_t) \subset |X_t|$，那么对 $t$ 的某个
邻域中的 $t' \in T$ 也有同一性质。
\end{proof}

\begin{lemma}
\label{moduli-lemma-quot-closed}
设 $B$ 是代数空间。设 $\pi : X \to Y$ 是代数空间之间的闭浸入，
且这些代数空间在 $B$ 上分离并有限表示。设 $\mathcal{F}$ 是拟凝聚
$\mathcal{O}_X$-模。则引理 \ref{moduli-lemma-quot-functorial} 的态射
$\Quotfunctor_{\mathcal{F}/X/B} \to
\Quotfunctor_{\pi_*\mathcal{F}/Y/B}$ 是同构。
\end{lemma}

\begin{proof}
对于每个 $B$ 上概形 $T$，态射 $\pi_T : X_T \to Y_T$ 是闭浸入。
于是，$\pi_{T, *}$ 给出 $\QCoh(\mathcal{O}_{X_T})$ 与
$\QCoh(\mathcal{O}_{Y_T})$ 的全子范畴之间的范畴等价；后者的对象
是被 $X_T$ 的理想层零化的拟凝聚模，见《空间的态射》中的引理
\ref{spaces-morphisms-lemma-i-star-equivalence}。
% P12-ERR-0202: see ../control/ERRATA_CANDIDATES.jsonl
由于 $(\pi_*\mathcal{F})_T$ 的商被该理想零化，便得到映射
$\Quotfunctor_{\mathcal{F}/X/B}(T) \to
\Quotfunctor_{\pi_*\mathcal{F}/Y/B}(T)$ 对所有 $T$ 都是双射，
这正是所需结论。
\end{proof}

\begin{lemma}
\label{moduli-lemma-quot-quotient}
设 $X \to B$ 如本节引言所述，并设
$\mathcal{F} \to \mathcal{G}$ 是拟凝聚 $\mathcal{O}_X$-模的满射。
则存在典范闭浸入
$\Quotfunctor_{\mathcal{G}/X/B} \to
\Quotfunctor_{\mathcal{F}/X/B}$。
\end{lemma}

\begin{proof}
令 $\mathcal{K} = \Ker(\mathcal{F} \to \mathcal{G})$。由拉回的
右正合性，对所有 $B$ 上概形 $T$，有
% P12-ERR-0203: see ../control/ERRATA_CANDIDATES.jsonl
$\mathcal{K}_T \to \mathcal{F}_T \to \mathcal{G}_T \to 0$
为正合列。特别地，$\mathcal{G}_T$ 的一个商确定
$\mathcal{F}_T$ 的一个商，从而得到函子间的变换
$\Quotfunctor_{\mathcal{G}/X/B} \to
\Quotfunctor_{\mathcal{F}/X/B}$。
由《空间上的平坦性》中的引理
\ref{spaces-flat-lemma-F-zero-closed-proper}，该变换是闭浸入。
% P12-ERR-0204: see ../control/ERRATA_CANDIDATES.jsonl
具体而言，给定 $\Quotfunctor_{\mathcal{F}/X/B}(T)$ 的元素
$\mathcal{F}_T \to \mathcal{Q}$，它向 $T'/T$ 的拉回位于该变换
的像中，当且仅当 $\mathcal{K}_{T'} \to \mathcal{Q}_{T'}$ 为零。
\end{proof}

\begin{remark}[数值不变量]
\label{moduli-remark-quot-numerical}
设 $f : X \to B$ 和 $\mathcal{F}$ 如本节引言所述。设 $I$ 是集合，
且对 $i \in I$，设 $E_i \in D(\mathcal{O}_X)$ 是完美对象。
设 $P : I \to \mathbf{Z}$ 是函数。回忆我们有态射
$$
\Quotfunctor_{\mathcal{F}/X/B} \longrightarrow \Cohstack_{X/B}
$$
它把 $\Quotfunctor_{\mathcal{F}/X/B}(T)$ 的元素
$\mathcal{F}_T \to \mathcal{Q}$ 送到 $\Cohstack_{X/B}$ 在 $T$
上的对象 $\mathcal{Q}$，见《Quot》中的命题
\ref{quot-proposition-quot} 的证明。因此可以作纤维积图
$$
\xymatrix{
\Quotfunctor^P_{\mathcal{F}/X/B} \ar[r] \ar[d] &
\Cohstack^P_{X/B} \ar[d] \\
\Quotfunctor_{\mathcal{F}/X/B} \ar[r] &
\Cohstack_{X/B}
}
$$
这就是左上角代数空间的定义图。左侧竖直箭头是平坦闭浸入；例如，
若 $I$ 有限，或 $B$ 局部 Noether，或 $I = \mathbf{Z}$ 且对某个
可逆 $\mathcal{O}_X$-模 $\mathcal{L}$ 有
$E_i = \mathcal{L}^{\otimes i}$，则它是既开又闭浸入（在最后一种
情形中，有时采用记号
$\Quotfunctor^{P, \mathcal{L}}_{\mathcal{F}/X/B}$）。
见情形 \ref{moduli-situation-numerical}、引理 \ref{moduli-lemma-open-P} 和
\ref{moduli-lemma-finite-list-perfect-objects}，以及例
\ref{moduli-example-hilbert-polynomial}。
\end{remark}

\begin{lemma}
\label{moduli-lemma-quot-tensor-invertible}
设 $f : X \to B$ 和 $\mathcal{F}$ 如本节引言所述。设
$\mathcal{L}$ 是可逆 $\mathcal{O}_X$-模。则与 $\mathcal{L}$
作张量积定义同构
$$
\Quotfunctor_{\mathcal{F}/X/B}
\longrightarrow
\Quotfunctor_{\mathcal{F} \otimes_{\mathcal{O}_X} \mathcal{L}/X/B}
$$
% P12-ERR-0205: see ../control/ERRATA_CANDIDATES.jsonl
给定数值多项式 $P(t)$，令 $P'(t) = P(t + 1)$，则此映射诱导既开又闭
子叠的同构
$\Quotfunctor^P_{\mathcal{F}/X/B}
\longrightarrow
\Quotfunctor^{P'}_{\mathcal{F} \otimes_{\mathcal{O}_X} \mathcal{L}/X/B}$。
\end{lemma}

\begin{proof}
令 $\mathcal{G} = \mathcal{F} \otimes_{\mathcal{O}_X} \mathcal{L}$。
注意
$\mathcal{G}_T = \mathcal{F}_T \otimes_{\mathcal{O}_{X_T}} \mathcal{L}_T$。
若 $\mathcal{F}_T \to \mathcal{Q}$ 是
$\Quotfunctor_{\mathcal{F}/X/B}(T)$ 的元素，就把它送到
$\Quotfunctor_{\mathcal{F} \otimes_{\mathcal{O}_X} \mathcal{L}/X/B}(T)$
的元素
$\mathcal{G}_T \to \mathcal{Q} \otimes_{\mathcal{O}_{X_T}} \mathcal{L}_T$。
这与拉回相容，因而按所需定义了函子间的变换。由于存在明显的逆变换，
它是同构。略去最后一个陈述的证明。
\end{proof}

\begin{lemma}
\label{moduli-lemma-quot-power-invertible}
设 $f : X \to B$ 和 $\mathcal{F}$ 如本节引言所述。设
$\mathcal{L}$ 是可逆 $\mathcal{O}_X$-模。则
$$
\Quotfunctor^{P, \mathcal{L}}_{\mathcal{F}/X/B} =
\Quotfunctor^{P', \mathcal{L}^{\otimes n}}_{\mathcal{F}/X/B}
$$
其中 $P'(t) = P(nt)$。
\end{lemma}

\begin{proof}
展开所有定义后立即得出。
\end{proof}






\section{Quot 的有界性}
\label{moduli-section-quot-bounded}

\noindent
与经典情形不同，我们已经知道 Quot 函子是代数空间，但尚不知道它是否
总能由有限型代数空间表示。

\begin{lemma}
\label{moduli-lemma-quot-Pn}
设 $n \geq 0$、$r \geq 1$、$P \in \mathbf{Q}[t]$。参数化
$\mathcal{O}_{\mathbf{P}^n_\mathbf{Z}}^{\oplus r}$ 的、Hilbert
多项式为 $P$ 的商的代数空间
$$
X = \Quotfunctor^P_{\mathcal{O}^{\oplus r}_{\mathbf{P}^n_\mathbf{Z}}/
\mathbf{P}^n_\mathbf{Z}/\mathbf{Z}}
$$
在 $\Spec(\mathbf{Z})$ 上真。
\end{lemma}

\begin{proof}
我们已经知道 $X \to \Spec(\mathbf{Z})$ 分离且局部有限表示
（引理 \ref{moduli-lemma-quot-s-lfp}）。我们还知道
$X \to \Spec(\mathbf{Z})$ 满足赋值判据的存在性部分，见引理
\ref{moduli-lemma-quot-existence-part}。由真性的赋值判据，只需证明我们的
Quot 空间拟紧，见《空间的态射》中的引理
\ref{spaces-morphisms-lemma-characterize-proper}。因此，只需找到一个
拟紧概形 $T$ 以及满态射 $T \to X$。设 $m$ 为《簇》中的引理
\ref{varieties-lemma-bound-quotients-free} 所给出的整数。令
$$
N = r{m + n \choose n} - P(m)
$$
以下把 $\mathbf{P}^n_\mathbf{Z} =
\text{Proj}(\mathbf{Z}[T_0, \ldots, T_n])$ 简记为
$\mathbf{P}^n$，未标底的积均指在 $\Spec(\mathbf{Z})$ 上的积。
证明的思路是在一个仿射概形 $T$ 上构造一个“万有”映射
$$
\Psi :
\mathcal{O}_{T \times \mathbf{P}^n}(-m)^{\oplus N}
\longrightarrow
\mathcal{O}_{T \times \mathbf{P}^n}^{\oplus r}
$$
并证明 $X$ 的每个点都对应于此映射在 $T$ 的某个点处的余核。

\medskip\noindent
$T$ 和 $\Psi$ 的定义。取 $T = \Spec(A)$，其中
$$
A = \mathbf{Z}[a_{i, j, E}]
$$
这里 $i \in \{1, \ldots, r\}$、$j \in \{1, \ldots, N\}$，
而 $E = (e_0, \ldots, e_n)$ 遍历总次数
% P12-ERR-0206: see ../control/ERRATA_CANDIDATES.jsonl
$|E| = \sum_{k = 0, \ldots n} e_k = m$ 的多重指标。然后把 $\Psi$
定义为一个矩阵，其 $(i, j)$ 元是映射
$$
\sum\nolimits_{E = (e_0, \ldots, e_n)}
a_{i, j, E} T_0^{e_0} \ldots T_n^{e_n} :
\mathcal{O}_{T \times \mathbf{P}^n}(-m)
\longrightarrow
\mathcal{O}_{T \times \mathbf{P}^n}
$$
其中求和遍历上述 $E$（当然，$i$ 和 $j$ 固定）。

\medskip\noindent
考虑 $T \times \mathbf{P}^n$ 上的商
$\mathcal{Q} = \Coker(\Psi)$。由《态射进阶》中的引理
\ref{more-morphisms-lemma-generic-flatness-stratification}，存在
$t \geq 0$ 以及闭子概形
$$
T = T_0 \supset T_1 \supset \ldots \supset T_t = \emptyset
$$
使得 $\mathcal{Q}$ 向
$(T_p \setminus T_{p + 1}) \times \mathbf{P}^n$ 的拉回
$\mathcal{Q}_p$ 在 $T_p \setminus T_{p + 1}$ 上平坦。注意，通过
拉回定义 $\mathcal{Q} = \Coker(\Psi)$ 的正合列，可得正合列
$$
\mathcal{O}_{(T_p \setminus T_{p + 1}) \times \mathbf{P}^n}(-m)^{\oplus N}
\to
\mathcal{O}_{(T_p \setminus T_{p + 1}) \times \mathbf{P}^n}^{\oplus r}
\to
\mathcal{Q}_p
\to
0
$$
因此得到态射
$$
% P12-ERR-0207: see ../control/ERRATA_CANDIDATES.jsonl
\coprod (T_p \setminus T_{p + 1})
\longrightarrow
% P12-ERR-0208: see ../control/ERRATA_CANDIDATES.jsonl
\Quotfunctor_{\mathcal{O}^{\oplus r}/\mathbf{P}/\mathbf{Z}}
\supset
\Quotfunctor^P_{\mathcal{O}^{\oplus r}/\mathbf{P}/\mathbf{Z}} = X
$$
由于左端是 Noether 概形，右端的包含是开的，只需证明 $X$ 的任意点
都在此态射的像中。

\medskip\noindent
设 $k$ 是域，并设 $x \in X(k)$。于是 $x$ 对应于凝聚
$\mathcal{O}_{\mathbf{P}^n_k}$-模的满射
$\mathcal{O}_{\mathbf{P}^n_k}^{\oplus r} \to \mathcal{F}$，
其中 $\mathcal{F}$ 的 Hilbert 多项式为 $P$。考虑短正合列
$$
0 \to \mathcal{K} \to
\mathcal{O}_{\mathbf{P}^n_k}^{\oplus r} \to
\mathcal{F} \to 0
$$
由《簇》中的引理 \ref{varieties-lemma-bound-quotients-free} 以及
$m$ 的选取，$\mathcal{K}$ 是 $m$-正则的。由《簇》中的引理
\ref{varieties-lemma-m-regular-globally-generated}，
$\mathcal{K}(m)$ 由整体截面生成。再由《簇》中的引理
\ref{varieties-lemma-m-regular-up} 和 $m$-正则性的定义，对
$i > 0$ 有
$H^i(\mathbf{P}^n_k, \mathcal{K}(m)) = 0$。因此
$$
\dim_k H^0(\mathbf{P}^n_k, \mathcal{K}(m)) =
\chi(\mathcal{K}(m)) =
\chi(\mathcal{O}_{\mathbf{P}^n_k}(m)^{\oplus r}) -
\chi(\mathcal{F}(m)) = N
$$
其中最后一个等式来自 $N$ 的选取。于是有满射
$$
\mathcal{O}_{\mathbf{P}^n_k}^{\oplus N}
\longrightarrow
\mathcal{K}(m)
$$
反向扭转并使用上述短正合列，可知 $\mathcal{F}$ 是映射
$$
\Psi_x :
\mathcal{O}_{\mathbf{P}^n_k}(-m)^{\oplus N}
\to
\mathcal{O}_{\mathbf{P}^n_k}^{\oplus r}
$$
的余核。存在唯一环同态 $\tau : A \to k$，使得 $\Psi$ 沿相应态射
$t = \Spec(\tau) : \Spec(k) \to T$ 的基变换等于 $\Psi_x$。
这是因为，定义 $\Psi_x$ 的
% P12-ERR-0209: see ../control/ERRATA_CANDIDATES.jsonl
$N \times r$ 矩阵的各项是 $T_0, \ldots, T_n$ 中次数为 $m$ 的
齐次多项式
$\sum \lambda_{i, j, E} T_0^{e_0} \ldots T_n^{e_n}$，其系数
$\lambda_{i, j, E} \in k$；我们可以令
$\tau(a_{i, j, E}) = \lambda_{i, j, E}$。于是对某个 $p$ 有
$t \in T_p \setminus T_{p + 1}$，且 $t$ 在上述态射下的像正是
$x$，如所需。
\end{proof}

\begin{lemma}
\label{moduli-lemma-quot-Pn-over-base}
设 $B$ 是代数空间，并令 $X = B \times \mathbf{P}^n_\mathbf{Z}$。
设 $\mathcal{L}$ 是 $\mathcal{O}_{\mathbf{P}^n}(1)$ 向 $X$ 的拉回。
设 $\mathcal{F}$ 是有限表示的 $\mathcal{O}_X$-模。参数化
$\mathcal{F}$ 的、关于 $\mathcal{L}$ 的 Hilbert 多项式为 $P$
的商的代数空间 $\Quotfunctor^P_{\mathcal{F}/X/B}$ 在 $B$ 上真。
\end{lemma}

\begin{proof}
此问题在 $B$ 上是 \'etale 局部的，见《空间的态射》中的引理
\ref{spaces-morphisms-lemma-proper-local}。因此可以假设 $B$ 是
仿射概形。此时 $\mathcal{L}$ 是 $X$ 上的丰富可逆模
（由《构造》中的引理 \ref{constructions-lemma-ample-on-proj}
以及《性质》中的定义 \ref{properties-definition-ample} 给出的
丰富可逆模之定义）。因此，由《性质》中的命题
\ref{properties-proposition-characterize-ample}，可以找到
$r' \geq 0$ 和 $r \geq 0$ 以及满射
$$
\mathcal{O}_X^{\oplus r} \longrightarrow
\mathcal{F} \otimes_{\mathcal{O}_X} \mathcal{L}^{\otimes r'}
$$
由引理 \ref{moduli-lemma-quot-tensor-invertible}，可以把 $\mathcal{F}$
替换为
$\mathcal{F} \otimes_{\mathcal{O}_X} \mathcal{L}^{\otimes r'}$，
并把 $P(t)$ 替换为 $P(t + r')$。由引理
\ref{moduli-lemma-quot-quotient}，得到闭浸入
$$
\Quotfunctor^P_{\mathcal{F}/X/B}
\longrightarrow
\Quotfunctor^P_{\mathcal{O}_X^{\oplus r}/X/B}
$$
由于我们已在引理 \ref{moduli-lemma-quot-Pn} 中证明
$\Quotfunctor^P_{\mathcal{O}_X^{\oplus r}/X/B} \to B$ 为真，
结论成立。
\end{proof}

\begin{lemma}
\label{moduli-lemma-quot-proper-over-base}
设 $f : X \to B$ 是代数空间之间的有限表示真态射。设
$\mathcal{F}$ 是有限表示的 $\mathcal{O}_X$-模。设
$\mathcal{L}$ 是在 $X/B$ 上丰富的可逆 $\mathcal{O}_X$-模，见
《空间上的除子》中的定义
\ref{spaces-divisors-definition-relatively-ample}。参数化
$\mathcal{F}$ 的、关于 $\mathcal{L}$ 的 Hilbert 多项式为 $P$
的商的代数空间 $\Quotfunctor^P_{\mathcal{F}/X/B}$ 在 $B$ 上真。
\end{lemma}

\begin{proof}
此问题在 $B$ 上是 \'etale 局部的，见《空间的态射》中的引理
\ref{spaces-morphisms-lemma-proper-local}。因此可以假设 $B$ 是
仿射概形。于是可以找到闭浸入 $i : X \to \mathbf{P}^n_B$，
使得对某个 $d \geq 1$ 有
$i^*\mathcal{O}_{\mathbf{P}^n_B}(1) \cong \mathcal{L}^{\otimes d}$。
见《态射》中的引理
\ref{morphisms-lemma-quasi-projective-finite-type-over-S}。
把 $\mathcal{L}$ 换成 $\mathcal{L}^{\otimes d}$，并把数值多项式
$P(t)$ 换成 $P(dt)$，不改变
$\Quotfunctor^P_{\mathcal{F}/X/B}$；略去一些细节。因此可以假设
$\mathcal{L} = i^*\mathcal{O}_{\mathbf{P}^n_B}(1)$。引理
\ref{moduli-lemma-quot-closed} 的同构
$\Quotfunctor_{\mathcal{F}/X/B} \to
\Quotfunctor_{i_*\mathcal{F}/\mathbf{P}^n_B/B}$ 诱导同构
$\Quotfunctor^P_{\mathcal{F}/X/B} \cong
\Quotfunctor^P_{i_*\mathcal{F}/\mathbf{P}^n_B/B}$。
由于由引理 \ref{moduli-lemma-quot-Pn-over-base}，
$\Quotfunctor^P_{i_*\mathcal{F}/\mathbf{P}^n_B/B}$ 在 $B$ 上真，
结论成立。
\end{proof}

\begin{lemma}
\label{moduli-lemma-quot-qc-over-base}
设 $f : X \to B$ 是代数空间之间的有限表示分离态射。设
$\mathcal{F}$ 是有限表示的 $\mathcal{O}_X$-模。设
$\mathcal{L}$ 是在 $X/B$ 上丰富的可逆 $\mathcal{O}_X$-模，见
《空间上的除子》中的定义
\ref{spaces-divisors-definition-relatively-ample}。参数化
$\mathcal{F}$ 的、关于 $\mathcal{L}$ 的 Hilbert 多项式为 $P$
的商的代数空间 $\Quotfunctor^P_{\mathcal{F}/X/B}$ 在 $B$
上分离且有限表示。
\end{lemma}

\begin{proof}
我们已经看到 $\Quotfunctor_{\mathcal{F}/X/B} \to B$ 分离且
局部有限表示，见引理 \ref{moduli-lemma-quot-s-lfp}。因此，只需证明注
\ref{moduli-remark-quot-numerical} 中的开子空间
$\Quotfunctor^P_{\mathcal{F}/X/B}$ 在 $B$ 上拟紧。

\medskip\noindent
此问题在 $B$ 上是 \'etale 局部的
（《空间的态射》中的引理
\ref{spaces-morphisms-lemma-quasi-compact-local}）。因此可以假设
$B$ 仿射。

\medskip\noindent
假设 $B = \Spec(\Lambda)$。把 $\Lambda = \colim \Lambda_i$
写成其有限型 $\mathbf{Z}$-子代数的余极限。于是可以找到一个 $i$，
以及引理中那样的、定义在 $B_i = \Spec(\Lambda_i)$ 上的系统
$X_i, \mathcal{F}_i, \mathcal{L}_i$，其向 $B$ 的基变换给出
$X, \mathcal{F}, \mathcal{L}$。这由《空间的极限》中的下列引理
得到：\ref{spaces-limits-lemma-descend-finite-presentation}
（用于找到 $X_i$）、
\ref{spaces-limits-lemma-descend-modules-finite-presentation}
（用于找到 $\mathcal{F}_i$）、
\ref{spaces-limits-lemma-descend-invertible-modules}
（用于找到 $\mathcal{L}_i$），以及
\ref{spaces-limits-lemma-descend-separated}
（用于使 $X_i$ 分离）。由于
$$
\Quotfunctor_{\mathcal{F}/X/B} = B \times_{B_i}
\Quotfunctor_{\mathcal{F}_i/X_i/B_i}
$$
并且 $\Quotfunctor^P_{\mathcal{F}/X/B}$ 也同样，问题约化为下一段
讨论的情形。

\medskip\noindent
假设 $B$ 仿射且 Noether。由引理
\ref{moduli-lemma-quot-power-invertible}，可以把 $\mathcal{L}$ 替换为
它的一个正幂。因此可以假设存在浸入
$i : X \to \mathbf{P}^n_B$，使得
$i^*\mathcal{O}_{\mathbf{P}^n}(1) = \mathcal{L}$。由《态射》中的
引理 \ref{morphisms-lemma-quasi-compact-immersion}，存在闭子概形
$X' \subset \mathbf{P}^n_B$，使得 $i$ 通过开浸入
$j : X \to X'$ 分解。由《性质》中的引理
\ref{properties-lemma-lift-finite-presentation}，存在有限表示的
$\mathcal{O}_{X'}$-模 $\mathcal{G}$，使得
$j^*\mathcal{G} = \mathcal{F}$。因此，由引理
\ref{moduli-lemma-quot-better-open} 得到开浸入
$$
\Quotfunctor_{\mathcal{F}/X/B}
\longrightarrow
\Quotfunctor_{\mathcal{G}/X'/B}
$$
显然，这个开浸入把 $\Quotfunctor^P_{\mathcal{F}/X/B}$ 送入
$\Quotfunctor^P_{\mathcal{G}/X'/B}$。现在，由引理
\ref{moduli-lemma-quot-proper-over-base}，
$\Quotfunctor^P_{\mathcal{G}/X'/B}$ 在 $B$ 上真。因此它是
Noether 的，而 Noether 代数空间的任意开子空间都拟紧，故结论成立。
\end{proof}




\section{Hilbert 函子的性质}
\label{moduli-section-hilb}

\noindent
设 $f : X \to B$ 是分离且有限表示的代数空间态射。则
$\Hilbfunctor_{X/B}$ 是在 $B$ 上局部有限表示的代数空间。
见《Quot》中的命题 \ref{quot-proposition-hilb}。

\begin{lemma}
\label{moduli-lemma-hilb-diagonal-closed}
$\Hilbfunctor_{X/B} \to B$ 的对角态射是有限表示的闭浸入。
\end{lemma}

\begin{proof}
我们已在《Quot》中的引理 \ref{quot-lemma-hilb-is-quot} 看到
$\Hilbfunctor_{X/B} = \Quotfunctor_{\mathcal{O}_X/X/B}$。
因此结论由引理 \ref{moduli-lemma-quot-diagonal-closed} 得出。
\end{proof}

\begin{lemma}
\label{moduli-lemma-hilb-s-lfp}
态射 $\Hilbfunctor_{X/B} \to B$ 分离且局部有限表示。
\end{lemma}

\begin{proof}
要检验 $\Hilbfunctor_{X/B} \to B$ 分离，需要证明其对角态射为
闭浸入。这由引理 \ref{moduli-lemma-hilb-diagonal-closed} 得出。
第二个陈述是《Quot》中的命题 \ref{quot-proposition-hilb} 的一部分。
\end{proof}

\begin{lemma}
\label{moduli-lemma-hilb-existence-part}
假设 $X \to B$ 为真且有限表示。则
$\Hilbfunctor_{X/B} \to B$ 满足赋值判据的存在性部分
（《空间的态射》中的定义
\ref{spaces-morphisms-definition-valuative-criterion}）。
\end{lemma}

\begin{proof}
我们已在《Quot》中的引理 \ref{quot-lemma-hilb-is-quot} 看到
$\Hilbfunctor_{X/B} = \Quotfunctor_{\mathcal{O}_X/X/B}$。
因此结论由引理 \ref{moduli-lemma-quot-existence-part} 得出。
\end{proof}

\begin{lemma}
\label{moduli-lemma-hilb-open}
设 $B$ 是代数空间。设 $\pi : X \to Y$ 是代数空间之间的开浸入，
且这些代数空间在 $B$ 上分离并有限表示。则 $\pi$ 诱导开浸入
$\Hilbfunctor_{X/B} \to \Hilbfunctor_{Y/B}$。
\end{lemma}

\begin{proof}
略。提示：若 $Z \subset X_T$ 是在 $T$ 上真的闭子概形，则 $Z$
也是 $Y_T$ 中的闭子概形。因此得到变换
$\Hilbfunctor_{X/B} \to \Hilbfunctor_{Y/B}$。若
$Z \subset Y_T$ 是 $\Hilbfunctor_{Y/B}(T)$ 的元素，并且对
$t \in T$ 有 $|Z_t| \subset |X_t|$，那么对 $t$ 的某个邻域中的
$t' \in T$ 也有同一性质。
\end{proof}

\begin{lemma}
\label{moduli-lemma-hilb-closed}
设 $B$ 是代数空间。设 $\pi : X \to Y$ 是代数空间之间的闭浸入，
且这些代数空间在 $B$ 上分离并有限表示。则 $\pi$ 诱导闭浸入
$\Hilbfunctor_{X/B} \to \Hilbfunctor_{Y/B}$。
\end{lemma}

\begin{proof}
由于 $\pi$ 是闭浸入，显然，给定闭子概形 $Z \subset X_T$，
% P12-ERR-0210: see ../control/ERRATA_CANDIDATES.jsonl
可把 $Z$ 看作 $X_T$ 的闭子概形。因此得到变换
$\Hilbfunctor_{X/B} \to \Hilbfunctor_{Y/B}$。立即可见该变换是
单态射。要证明它是闭浸入，可以对映射
$\mathcal{O}_Y \to \mathcal{O}_X$ 使用引理
\ref{moduli-lemma-quot-quotient}，并使用《Quot》中的引理
\ref{quot-lemma-hilb-is-quot} 的等同
$\Hilbfunctor_{X/B} = \Quotfunctor_{\mathcal{O}_X/X/B}$、
$\Hilbfunctor_{Y/B} = \Quotfunctor_{\mathcal{O}_Y/Y/B}$。
\end{proof}

\begin{remark}[数值不变量]
\label{moduli-remark-hilb-numerical}
设 $f : X \to B$ 如本节引言所述。设 $I$ 是集合，且对 $i \in I$，
设 $E_i \in D(\mathcal{O}_X)$ 是完美对象。设
$P : I \to \mathbf{Z}$ 是函数。回忆
$\Hilbfunctor_{X/B} = \Quotfunctor_{\mathcal{O}_X/X/B}$，见
《Quot》中的引理 \ref{quot-lemma-hilb-is-quot}。因此可以定义
$$
\Hilbfunctor^P_{X/B} = \Quotfunctor^P_{\mathcal{O}_X/X/B}
$$
其中 $\Quotfunctor^P_{\mathcal{O}_X/X/B}$ 如注
\ref{moduli-remark-quot-numerical} 所述。态射
$$
\Hilbfunctor^P_{X/B} \longrightarrow \Hilbfunctor_{X/B}
$$
是平坦闭浸入；例如，若 $I$ 有限，或 $B$ 局部 Noether，或
$I = \mathbf{Z}$ 且对某个可逆 $\mathcal{O}_X$-模 $\mathcal{L}$
有 $E_i = \mathcal{L}^{\otimes i}$，则它是既开又闭浸入。
在最后一种情形中，有时采用记号
$\Hilbfunctor^{P, \mathcal{L}}_{X/B}$。
\end{remark}

\begin{lemma}
\label{moduli-lemma-hilb-proper-over-base}
设 $f : X \to B$ 是代数空间之间的有限表示真态射。设
$\mathcal{L}$ 是在 $X/B$ 上丰富的可逆 $\mathcal{O}_X$-模，见
《空间上的除子》中的定义
\ref{spaces-divisors-definition-relatively-ample}。参数化关于
$\mathcal{L}$ 的 Hilbert 多项式为 $P$ 的闭子概形的代数空间
$\Hilbfunctor^P_{X/B}$ 在 $B$ 上真。
\end{lemma}

\begin{proof}
回忆 $\Hilbfunctor_{X/B} =
\Quotfunctor_{\mathcal{O}_X/X/B}$，见《Quot》中的引理
\ref{quot-lemma-hilb-is-quot}。因此本引理立即由引理
\ref{moduli-lemma-quot-proper-over-base} 得出。
\end{proof}

\begin{lemma}
\label{moduli-lemma-hilb-qc-over-base}
设 $f : X \to B$ 是代数空间之间的有限表示分离态射。设
$\mathcal{L}$ 是在 $X/B$ 上丰富的可逆 $\mathcal{O}_X$-模，见
《空间上的除子》中的定义
\ref{spaces-divisors-definition-relatively-ample}。参数化关于
$\mathcal{L}$ 的 Hilbert 多项式为 $P$ 的闭子概形的代数空间
$\Hilbfunctor^P_{X/B}$ 在 $B$ 上分离且有限表示。
\end{lemma}

\begin{proof}
回忆 $\Hilbfunctor_{X/B} =
\Quotfunctor_{\mathcal{O}_X/X/B}$，见《Quot》中的引理
\ref{quot-lemma-hilb-is-quot}。因此本引理立即由引理
\ref{moduli-lemma-quot-qc-over-base} 得出。
\end{proof}









\section{Picard 叠的性质}
\label{moduli-section-picard-stack}

\noindent
设 $f : X \to B$ 是平坦、真且有限表示的代数空间态射。则参数化
$X/B$ 上可逆层的叠 $\Picardstack_{X/B}$ 是代数的，见《Quot》中的
命题 \ref{quot-proposition-pic}。

\begin{lemma}
\label{moduli-lemma-pic-diagonal-affine-fp}
$\Picardstack_{X/B}$ 在 $B$ 上的对角态射仿射且有限表示。
\end{lemma}

\begin{proof}
我们已在《Quot》中的引理
\ref{quot-lemma-picard-stack-open-in-coh} 看到，
$\Picardstack_{X/B}$ 是 $\Cohstack_{X/B}$ 的开子叠。
因此结论由引理 \ref{moduli-lemma-coherent-diagonal-affine-fp} 得出。
\end{proof}

\begin{lemma}
\label{moduli-lemma-pic-qs-lfp}
态射 $\Picardstack_{X/B} \to B$ 拟分离且局部有限表示。
\end{lemma}

\begin{proof}
我们已在《Quot》中的引理
\ref{quot-lemma-picard-stack-open-in-coh} 看到，
$\Picardstack_{X/B}$ 是 $\Cohstack_{X/B}$ 的开子叠。
因此结论由引理 \ref{moduli-lemma-coherent-qs-lfp} 得出。
\end{proof}

\begin{lemma}
\label{moduli-lemma-pic-existence-part}
除为真之外，再假设 $X \to B$ 光滑。则
$\Picardstack_{X/B} \to B$ 满足赋值判据的存在性部分
（《叠的态射》中的定义
\ref{stacks-morphisms-definition-existence}）。
\end{lemma}

\begin{proof}
作基变换后，问题立即约化为下述问题：给定分式域为 $K$ 的赋值环
$R$、一个在 $R$ 上真且光滑的代数空间 $X$，以及一个可逆
$\mathcal{O}_{X_K}$-模 $\mathcal{L}_K$，证明存在一个可逆
$\mathcal{O}_X$-模 $\mathcal{L}$，其泛纤维为 $\mathcal{L}_K$。
注意，$X_K$ 是 Noether、分离且正则的
（使用《空间的态射》中的引理
\ref{spaces-morphisms-lemma-finite-presentation-noetherian}
和《域上的空间》中的引理
\ref{spaces-over-fields-lemma-smooth-regular}）。因此，对
$X_K$ 中的两个有效 Cartier 除子 $D_K, D'_K$，可以在 Picard 群中把
$\mathcal{L}_K$ 写成 $\mathcal{O}_{X_K}(D_K)$ 与
$\mathcal{O}_{X_K}(D'_K)$ 之差，见《空间上的除子》中的引理
\ref{spaces-divisors-lemma-Noetherian-regular-separated-pic-effective-Cartier}。
最后，我们知道 $D_K$ 和 $D'_K$ 分别是有效 Cartier 除子
$D, D' \subset X$ 的限制，见《空间上的除子》中的引理
\ref{spaces-divisors-lemma-smooth-over-valuation-ring-effective-Cartier}。
\end{proof}

\begin{lemma}
\label{moduli-lemma-pic-inertia}
假设对所有 $B$ 上概形 $T$ 都有
$f_{T, *}\mathcal{O}_{X_T} \cong \mathcal{O}_T$。则
$\Picardstack_{X/B}$ 的惯性叠等于
$\mathbf{G}_m \times \Picardstack_{X/B}$。
\end{lemma}

\begin{proof}
这在《叠的例子》中的例
\ref{examples-stacks-example-inertia-stack-of-picard} 中已有说明。
\end{proof}

\begin{lemma}
\label{moduli-lemma-pic-curves-smooth}
除本节的其他假设外，再假设 $f : X \to B$ 的相对维数
$\leq 1$。则 $\Picardstack_{X/B} \to B$ 光滑。
\end{lemma}

\begin{proof}
由引理 \ref{moduli-lemma-pic-qs-lfp}，我们已知
$\Picardstack_{X/B} \to B$ 局部有限表示。因此，由《叠的态射进阶》
中的引理 \ref{stacks-more-morphisms-lemma-smooth-formally-smooth}，
只需证明 $\Picardstack_{X/B} \to B$ 形式光滑。
作基变换后，问题立即约化为下述问题：给定仿射概形的一阶加厚
$T \subset T'$，给定真、平坦、有限表示且相对维数 $\leq 1$ 的
$X' \to T'$，并对 $X = T \times_{T'} X'$ 给定一个可逆
$\mathcal{O}_X$-模 $\mathcal{L}$，证明存在一个可逆
$\mathcal{O}_{X'}$-模 $\mathcal{L}'$，其在 $X$ 上的限制为
$\mathcal{L}$。由于 $T \subset T'$ 是一阶加厚，
$X \subset X'$ 也同样，见《空间的态射进阶》中的引理
\ref{spaces-more-morphisms-lemma-base-change-thickening}。
由《空间的态射进阶》中的引理
\ref{spaces-more-morphisms-lemma-picard-group-first-order-thickening}，
只需证明 $H^2(X, \mathcal{I}) = 0$，其中 $\mathcal{I}$ 是在
$X'$ 中截出 $X$ 的拟凝聚理想。
% P12-ERR-0211: see ../control/ERRATA_CANDIDATES.jsonl
以 $f : X \to T$ 表示结构态射。由《空间的上同调》中的引理
\ref{spaces-cohomology-lemma-higher-direct-images-zero-above-dimension-fibre}，
对 $p > 1$ 有 $R^pf_*\mathcal{I} = 0$。因此，由《空间的上同调》
中的引理
\ref{spaces-cohomology-lemma-quasi-coherence-higher-direct-images-application}
得到所需消失（这里最终用到了 $T$ 仿射）。
\end{proof}




\section{Picard 函子的性质}
\label{moduli-section-picard-functor}

\noindent
设 $f : X \to B$ 是平坦、真且有限表示的代数空间态射，并且还假设
对每个 $T/B$，典范映射
$$
\mathcal{O}_T \longrightarrow f_{T, *}\mathcal{O}_{X_T}
$$
都是同构。则 Picard 函子 $\Picardfunctor_{X/B}$ 是代数空间，
见《Quot》中的命题 \ref{quot-proposition-pic-functor}。
% P12-ERR-0212: see ../control/ERRATA_CANDIDATES.jsonl
它与 Picard 叠有紧密关系。

\begin{lemma}
\label{moduli-lemma-pic-gerbe-over-pic-functor}
态射 $\Picardstack_{X/B} \to \Picardfunctor_{X/B}$ 使 Picard 叠
成为 Picard 函子上的 gerbe。
\end{lemma}

\begin{proof}
$\Picardstack_{X/B} \to \Picardfunctor_{X/B}$ 为 gerbe 的定义
见《叠的态射》中的定义
\ref{stacks-morphisms-definition-gerbe}；该定义又引用《叠》中的定义
\ref{stacks-definition-gerbe-over-stack-in-groupoids}。为证明此事，
我们检验《叠》中的引理 \ref{stacks-lemma-when-gerbe} 的条件
(2)(a) 和 (2)(b)。这立即由《Quot》中的引理
\ref{quot-lemma-pic-over-pic} 得出；下面给出详细说明。

\medskip\noindent
条件 (2)(a)。设对某个 $B$ 上概形 $U$ 有
$\xi \in \Picardfunctor_{X/B}(U)$。由于
$\Picardfunctor_{X/B}$ 是 $B$ 上概形上的规则
$T \mapsto \Pic(X_T)$ 的 fppf 层化
（《Quot》中的情形 \ref{quot-situation-pic}），存在一个 fppf 覆盖
$\{U_i \to U\}$，使得 $\xi|_{U_i}$ 对应于 $X_{U_i}$ 上的某个
可逆模 $\mathcal{L}_i$。于是
$(U_i \to B, \mathcal{L}_i)$ 是 $\Picardstack_{X/B}$ 在 $U_i$
上的对象，并映到 $\xi|_{U_i}$。

\medskip\noindent
条件 (2)(b)。设 $U$ 是 $B$ 上概形，并设 $\mathcal{L}, \mathcal{N}$
是 $X_U$ 上的可逆模，它们映到 $\Picardfunctor_{X/B}(U)$ 的同一
元素。则存在 fppf 覆盖 $\{U_i \to U\}$，使得
$\mathcal{L}|_{X_{U_i}}$ 同构于
$\mathcal{N}|_{X_{U_i}}$。因此按所需找到同构
$(U \to B, \mathcal{L})|_{U_i} \to
(U \to B, \mathcal{N})|_{U_i}$。
\end{proof}

\begin{lemma}
\label{moduli-lemma-pic-functor-diagonal-qc-immersion}
$\Picardfunctor_{X/B}$ 在 $B$ 上的对角态射是拟紧浸入。
\end{lemma}

\begin{proof}
由《Quot》中的引理 \ref{quot-lemma-diagonal-pic}，对角态射是浸入。
为完成证明，只需说明它拟紧。由引理
\ref{moduli-lemma-pic-diagonal-affine-fp}，
$\Picardstack_{X/B}$ 的对角态射拟紧；由引理
\ref{moduli-lemma-pic-gerbe-over-pic-functor}，
$\Picardstack_{X/B}$ 是 $\Picardfunctor_{X/B}$ 上的 gerbe。
结论由《叠的态射》中的引理
\ref{stacks-morphisms-lemma-gerbe-diagonal-quasi-compact} 得出。
\end{proof}

\begin{lemma}
\label{moduli-lemma-pic-functor-qs-lfp}
态射 $\Picardfunctor_{X/B} \to B$ 拟分离且局部有限表示。
\end{lemma}

\begin{proof}
要检验 $\Picardfunctor_{X/B} \to B$ 拟分离，需要证明其对角态射
拟紧。这立即由引理
\ref{moduli-lemma-pic-functor-diagonal-qc-immersion} 得出。
态射 $\Picardstack_{X/B} \to \Picardfunctor_{X/B}$ 是满、平坦
且局部有限表示的
（由引理 \ref{moduli-lemma-pic-gerbe-over-pic-functor} 和《叠的态射》中的
引理 \ref{stacks-morphisms-lemma-gerbe-fppf}）。因此，由《叠的态射》
中的引理
\ref{stacks-morphisms-lemma-flat-finite-presentation-permanence}，
只需证明 $\Picardstack_{X/B} \to B$ 局部有限表示。
这由引理 \ref{moduli-lemma-pic-qs-lfp} 得出。
\end{proof}

\begin{lemma}
\label{moduli-lemma-pic-functor-uniqueness-part}
除本节的其他假设外，再假设 $X \to B$ 的几何纤维整。则
$\Picardfunctor_{X/B} \to B$ 分离。
\end{lemma}

\begin{proof}
由于 $\Picardfunctor_{X/B} \to B$ 拟分离，只需检验赋值判据的
唯一性部分，见《空间的态射》中的引理
\ref{spaces-morphisms-lemma-valuative-criterion-separatedness}。
这立即约化为下述问题：给定
\begin{enumerate}
\item 一个分式域为 $K$ 的赋值环 $R$；
% P12-ERR-0213: see ../control/ERRATA_CANDIDATES.jsonl
\item 一个在 $R$ 上真且平坦、具有整几何纤维的代数空间 $X$；
\item 一个元素 $a \in \Picardfunctor_{X/R}(R)$，满足
$a|_{\Spec(K)} = 0$，
\end{enumerate}
需要证明 $a = 0$。把《叠的态射》中的引理
\ref{stacks-morphisms-lemma-lift-valuation-ring-through-flat-morphism}
应用于满平坦态射
$\Picardstack_{X/R} \to \Picardfunctor_{X/R}$
（由引理 \ref{moduli-lemma-pic-gerbe-over-pic-functor} 和《叠的态射》中的
引理 \ref{stacks-morphisms-lemma-gerbe-fppf}，它是满且平坦的），
在把 $R$ 替换为一个扩张后，可以假设 $a$ 由一个可逆
$\mathcal{O}_X$-模 $\mathcal{L}$ 给出。由于
$a|_{\Spec(K)} = 0$，由《Quot》中的引理
\ref{quot-lemma-flat-geometrically-connected-fibres} 有
$\mathcal{L}_K \cong \mathcal{O}_{X_K}$。

\medskip\noindent
% P12-ERR-0214: see ../control/ERRATA_CANDIDATES.jsonl
以 $f : X \to \Spec(R)$ 表示结构态射。设
$\eta, 0 \in \Spec(R)$ 分别为泛点与闭点。考虑
$\Spec(R)$ 上的完美复形
$K = Rf_*\mathcal{L}$ 和
$M = Rf_*(\mathcal{L}^{\otimes -1})$，见《空间的导出范畴》中的
引理
\ref{spaces-perfect-lemma-flat-proper-perfect-direct-image-general}。
考虑《空间的导出范畴》中的引理
\ref{spaces-perfect-lemma-jump-loci} 与 $K$ 和 $M$ 关联的函数
$\beta_{K, i}, \beta_{M, i} : \Spec(R) \to \mathbf{Z}$。
% P12-ERR-0215: see ../control/ERRATA_CANDIDATES.jsonl
由于 $K$ 与 $M$ 的形成与基变换交换（见上述引理），由《域上的空间》
中的引理
\ref{spaces-over-fields-lemma-proper-geometrically-reduced-global-sections}
以及关于 $f$ 的纤维的假设，有
$\beta_{K, 0}(\eta) = \beta_{M, 0}(\eta) = 1$。
由上半连续性，
$\beta_{K, 0}(0) \geq 1$ 且 $\beta_{M, 0}(0) \geq 1$。
由《域上的空间》中的引理
\ref{spaces-over-fields-lemma-characterize-trivial-pic-integral}，
可知 $\mathcal{L}$ 在特殊纤维 $X_0$ 上的限制平凡。于是如上得到
% P12-ERR-0216: see ../control/ERRATA_CANDIDATES.jsonl
$\beta_{K, 0}(0) = \beta_{M, 0} = 1$。
然后由《代数进阶》中的引理
\ref{more-algebra-lemma-lift-pseudo-coherent-from-residue-field}，
可以用下列形式的复形表示 $K$：
$$
\ldots \to 0 \to R \to R^{\oplus \beta_{K, 1}(0)} \to
R^{\oplus \beta_{K, 2}(0)} \to \ldots
$$
现在 $R \to R^{\oplus \beta_{K, 1}(0)}$ 为零，因为
$\beta_{K, 0}(\eta) = 1$。换言之，在 $D(R)$ 中
$K = R \oplus \tau_{\geq 1}(K)$，其中对某个
$b \in \mathbf{Z}$，$\tau_{\geq 1}(K)$ 的 Tor 振幅位于
$[1, b]$。因此存在整体截面 $s \in H^0(X, \mathcal{L})$，
其在 $X_0$ 上的限制 $s_0$ 处处非零
（仍因为 $K$ 的形成与基变换交换）。于是
$s : \mathcal{O}_X \to \mathcal{L}$ 是可逆层之间的映射，其在
$X_0$ 上的限制是同构，因而它本身也是同构，如所需。
\end{proof}

\begin{lemma}
\label{moduli-lemma-pic-functor-curves-smooth}
除本节的其他假设外，再假设 $f : X \to B$ 的相对维数
$\leq 1$。则 $\Picardfunctor_{X/B} \to B$ 光滑。
\end{lemma}

\begin{proof}
由引理 \ref{moduli-lemma-pic-curves-smooth}，我们知道
$\Picardstack_{X/B} \to B$ 光滑。结合引理
\ref{moduli-lemma-pic-gerbe-over-pic-functor} 与《叠的态射》中的引理
\ref{stacks-morphisms-lemma-gerbe-smooth}，态射
$\Picardstack_{X/B} \to \Picardfunctor_{X/B}$ 满且光滑。
因此，若 $U$ 是概形且
$U \to \Picardstack_{X/B}$ 满且光滑，则
$U \to \Picardfunctor_{X/B}$ 满且光滑，而 $U \to B$ 也满且
光滑（因为这些性质在复合下保持）。所以，例如由《空间上的下降》
中的引理
\ref{spaces-descent-lemma-syntomic-smooth-etale-permanence}，
$\Picardfunctor_{X/B} \to B$ 光滑。
\end{proof}






\section{相对态射的性质}
\label{moduli-section-relative-morphisms}

\noindent
设 $B$ 是代数空间。设 $X$ 和 $Y$ 是 $B$ 上代数空间，其中
$Y \to B$ 平坦、真且有限表示，而 $X \to B$ 分离且有限表示。
则相对态射的函子 $\mathit{Mor}_B(Y, X)$ 是在 $B$ 上局部有限表示
的代数空间。见《Quot》中的命题 \ref{quot-proposition-Mor}。

\begin{lemma}
\label{moduli-lemma-Mor-diagonal-closed}
$\mathit{Mor}_B(Y, X) \to B$ 的对角态射是有限表示的闭浸入。
\end{lemma}

\begin{proof}
存在开浸入
$\mathit{Mor}_B(Y, X) \to \Hilbfunctor_{Y \times_B X/B}$，见
《Quot》中的引理 \ref{quot-lemma-Mor-into-Hilb-open}。
因此本引理由引理 \ref{moduli-lemma-hilb-diagonal-closed} 得出。
\end{proof}

\begin{lemma}
\label{moduli-lemma-Mor-s-lfp}
态射 $\mathit{Mor}_B(Y, X) \to B$ 分离且局部有限表示。
\end{lemma}

\begin{proof}
要检验 $\mathit{Mor}_B(Y, X) \to B$ 分离，需要证明其对角态射为
闭浸入。这由引理 \ref{moduli-lemma-Mor-diagonal-closed} 得出。
第二个陈述是《Quot》中的命题 \ref{quot-proposition-Mor} 的一部分。
\end{proof}

\begin{lemma}
\label{moduli-lemma-Isom-in-Mor}
设 $B, X, Y$ 如本节引言所述，并额外假设 $X \to B$ 为真。
则由同构组成的子函子
$\mathit{Isom}_B(Y, X) \subset \mathit{Mor}_B(Y, X)$
是开子空间。
\end{lemma}

\begin{proof}
这立即由《空间的态射进阶》中的引理
\ref{spaces-more-morphisms-lemma-where-isomorphism} 得出。
\end{proof}

\begin{remark}[数值不变量]
\label{moduli-remark-Mor-numerical}
设 $B, X, Y$ 如本节引言所述。设 $I$ 是集合，且对 $i \in I$，
设 $E_i \in D(\mathcal{O}_{Y \times_B X})$ 是完美对象。
设 $P : I \to \mathbf{Z}$ 是函数。回忆
$$
\mathit{Mor}_B(Y, X) \subset
\Hilbfunctor_{Y \times_B X/B}
$$
是开子空间，见《Quot》中的引理
\ref{quot-lemma-Mor-into-Hilb-open}。因此可以定义
$$
\mathit{Mor}^P_B(Y, X) =
\mathit{Mor}_B(Y, X) \cap \Hilbfunctor^P_{Y \times_B X/B}
$$
其中 $\Hilbfunctor^P_{Y \times_B X/B}$ 如注
\ref{moduli-remark-hilb-numerical} 所述。态射
$$
\mathit{Mor}^P_B(Y, X) \longrightarrow \mathit{Mor}_B(Y, X)
$$
是平坦闭浸入；例如，若 $I$ 有限，或 $B$ 局部 Noether，或
$I = \mathbf{Z}$ 且对某个可逆
$\mathcal{O}_{Y \times_B X}$-模 $\mathcal{L}$ 有
$E_i = \mathcal{L}^{\otimes i}$，则它是既开又闭浸入。
在最后一种情形中，有时采用记号
$\mathit{Mor}^{P, \mathcal{L}}_B(Y, X)$。
\end{remark}

\begin{lemma}
\label{moduli-lemma-Mor-qc-over-base}
设 $B, X, Y$ 如本节引言所述，并设 $\mathcal{L}$ 在 $X/B$ 上丰富，
$\mathcal{N}$ 在 $Y/B$ 上丰富。见《空间上的除子》中的定义
\ref{spaces-divisors-definition-relatively-ample}。
设 $P$ 是数值多项式。则
$$
\mathit{Mor}^{P, \mathcal{M}}_B(Y, X) \longrightarrow B
$$
分离且有限表示，其中
$\mathcal{M} = \text{pr}_1^*\mathcal{N}
\otimes_{\mathcal{O}_{Y \times_B X}} \text{pr}_2^*\mathcal{L}$。
\end{lemma}

\begin{proof}
由引理 \ref{moduli-lemma-Mor-s-lfp}，态射
$\mathit{Mor}_B(Y, X) \to B$ 分离且局部有限表示。因此只需证明
注 \ref{moduli-remark-Mor-numerical} 的既开又闭子空间
$\mathit{Mor}^{P, \mathcal{M}}_B(Y, X)$ 在 $B$ 上拟紧。

\medskip\noindent
此问题在 $B$ 上是 \'etale 局部的
（《空间的态射》中的引理
\ref{spaces-morphisms-lemma-quasi-compact-local}）。
因此可以假设 $B$ 仿射。

\medskip\noindent
假设 $B = \Spec(\Lambda)$。注意，$X$ 和 $Y$ 是概形，而
$\mathcal{L}$ 和 $\mathcal{N}$ 分别是 $X$ 和 $Y$ 上的丰富可逆层
（这立即由定义得出）。把 $\Lambda = \colim \Lambda_i$ 写成其
有限型 $\mathbf{Z}$-子代数的余极限。于是可以找到一个 $i$，以及
引理中那样的、定义在 $B_i = \Spec(\Lambda_i)$ 上的系统
$X_i, Y_i, \mathcal{L}_i, \mathcal{N}_i$，其向 $B$ 的基变换给出
$X, Y, \mathcal{L}, \mathcal{N}$。这由《极限》中的下列引理得到：
\ref{limits-lemma-descend-finite-presentation}
（用于找到 $X_i$、$Y_i$）、
\ref{limits-lemma-descend-invertible-modules}
（用于找到 $\mathcal{L}_i$、$\mathcal{N}_i$）、
\ref{limits-lemma-descend-separated-finite-presentation}
（用于使 $X_i \to B_i$ 分离）、
\ref{limits-lemma-eventually-proper}
（用于使 $Y_i \to B_i$ 为真），以及
\ref{limits-lemma-limit-ample}
（用于使 $\mathcal{L}_i$、$\mathcal{N}_i$ 丰富）。由于
$$
\mathit{Mor}_B(Y, X) = B \times_{B_i} \mathit{Mor}_{B_i}(Y_i, X_i)
$$
% P12-ERR-0217: see ../control/ERRATA_CANDIDATES.jsonl
并且 $\mathit{Mor}^P_B(Y, X)$ 也同样，问题约化为下一段讨论的情形。

\medskip\noindent
假设 $B$ 是 Noether 仿射概形。由《性质》中的引理
\ref{properties-lemma-ample-on-product}，$\mathcal{M}$ 丰富。
由引理 \ref{moduli-lemma-hilb-qc-over-base}，
$\Hilbfunctor^{P, \mathcal{M}}_{Y \times_B X/B}$ 在 $B$ 上
有限表示，因而是 Noether 的。按构造，
$$
\mathit{Mor}^{P, \mathcal{M}}_B(Y, X) =
\mathit{Mor}_B(Y, X) \cap
\Hilbfunctor^{P, \mathcal{M}}_{Y \times_B X/B}
$$
是 $\Hilbfunctor^{P, \mathcal{M}}_{Y \times_B X/B}$ 的开子空间，
因而拟紧（因为 Noether 代数空间的开子空间拟紧）。
\end{proof}






\section{极化真概形叠的性质}
\label{moduli-section-polarized}

\noindent
本节讨论模叠
$$
\Polarizedstack \longrightarrow \Spec(\mathbf{Z})
$$
的性质；它在概形 $S$ 上的截面范畴，是由配备相对丰富可逆层的、
% P12-ERR-0218: see ../control/ERRATA_CANDIDATES.jsonl
在 $S$ 上真、平坦且有限表示的概形组成的范畴。由《Quot》中的定理
\ref{quot-theorem-polarized-algebraic}，这是代数叠。

\begin{lemma}
\label{moduli-lemma-polarized-diagonal-separated-fp}
$\Polarizedstack$ 的对角态射分离且有限表示。
\end{lemma}

\begin{proof}
回忆 $\Polarizedstack$ 是保持极限的代数叠，见《Quot》中的引理
\ref{quot-lemma-polarized-limits}。由《叠的极限》中的引理
\ref{stacks-limits-lemma-limit-preserving-diagonal}，这蕴含
$\Delta : \Polarizedstack \to
\Polarizedstack \times \Polarizedstack$ 保持极限。因此，由
《叠的极限》中的命题
\ref{stacks-limits-proposition-characterize-locally-finite-presentation}，
$\Delta$ 局部有限表示。

\medskip\noindent
下面证明 $\Delta$ 分离。为此，只需证明：给定仿射概形 $U$ 以及
$\Polarizedstack$ 在 $U$ 上的两个对象
$\upsilon = (Y, \mathcal{N})$ 和
$\chi = (X, \mathcal{L})$，代数空间
$$
\mathit{Isom}_{\Polarizedstack}(\upsilon, \chi)
$$
分离。把同构 $\upsilon_T \to \chi_T$ 送到其底层同构
$Y_T \to X_T$ 的规则定义态射
$$
\mathit{Isom}_{\Polarizedstack}(\upsilon, \chi)
\longrightarrow
\mathit{Isom}_U(Y, X)
$$
由引理 \ref{moduli-lemma-Mor-s-lfp} 和 \ref{moduli-lemma-Isom-in-Mor}，
我们已经看到其目标是分离代数空间，所以只需证明此态射分离。
% P12-ERR-0219: see ../control/ERRATA_CANDIDATES.jsonl
给定某个概形 $T/U$ 上的同构 $f : Y_T \to X_T$，显然有
$$
\mathit{Isom}_{\Polarizedstack}(\upsilon, \chi)
\times_{\mathit{Isom}_U(Y, X), [f]} T
=
\mathit{Isom}(\mathcal{N}_T, f^*\mathcal{L}_T)
$$
这里 $[f] : T \to \mathit{Isom}_U(Y, X)$ 表示与 $f$ 对应的
$T$-值点，而
$\mathit{Isom}(\mathcal{N}_T, f^*\mathcal{L}_T)$ 是章节
\ref{moduli-section-hom-isom} 中讨论的代数空间。
% P12-ERR-0220: see ../control/ERRATA_CANDIDATES.jsonl
由于该代数空间在 $U$ 上仿射，所断言的结论蕴含 $\Delta$ 分离。

\medskip\noindent
最后证明 $\Delta$ 拟紧。由于 $\Delta$ 可由代数空间表示，
% P12-ERR-0221: see ../control/ERRATA_CANDIDATES.jsonl
只需检验 $\Delta$ 沿一个满光滑态射
$U \to \Polarizedstack \times \Polarizedstack$ 的基变换拟紧
（例如见《叠的性质》中的引理
\ref{stacks-properties-lemma-check-property-covering}）。
可以假设 $U = \coprod U_i$ 是仿射开集的不交并。
由于 $\Polarizedstack$ 保持极限（见上文），
$\Polarizedstack \to \Spec(\mathbf{Z})$ 局部有限表示，因而
$U_i \to \Spec(\mathbf{Z})$ 局部有限表示
（《叠的极限》中的命题
\ref{stacks-limits-proposition-characterize-locally-finite-presentation}
以及《叠的态射》中的引理
\ref{stacks-morphisms-lemma-composition-finite-presentation} 和
\ref{stacks-morphisms-lemma-smooth-locally-finite-presentation}）。
特别地，$U_i$ 是 Noether 仿射概形。于是问题约化为下一段讨论的情形。

\medskip\noindent
在本段中，给定 Noether 仿射概形 $U$ 以及
$\Polarizedstack$ 在 $U$ 上的两个对象
$\upsilon = (Y, \mathcal{N})$ 和
$\chi = (X, \mathcal{L})$，我们证明代数空间
$$
\mathit{Isom}_{\Polarizedstack}(\upsilon, \chi)
$$
拟紧。由于 $U$ 的连通分支既开又闭，可以用这些分支替换 $U$。
因此可以并且确实假设 $U$ 连通。设 $u \in U$ 是一个点。设 $P$
为 Hilbert 多项式
$n \mapsto \chi(Y_u, \mathcal{N}_u^{\otimes n})$，见《簇》中的
引理 \ref{varieties-lemma-numerical-polynomial-from-euler}。
由于 $U$ 连通，并且函数
$u \mapsto \chi(Y_u, \mathcal{N}_u^{\otimes n})$ 局部常值
（见《概形的导出范畴》中的引理
\ref{perfect-lemma-chi-locally-constant-geometric}），可知在 $U$ 的
每个点都得到同一 Hilbert 多项式。在 $Y \times_U X$ 上令
$\mathcal{M} = \text{pr}_1^*\mathcal{N}
\otimes_{\mathcal{O}_{Y \times_U X}} \text{pr}_2^*\mathcal{L}$。
% P12-ERR-0222: see ../control/ERRATA_CANDIDATES.jsonl
给定某个 $U$ 上概形 $T$ 的元素
$(f, \varphi) \in
\mathit{Isom}_{\Polarizedstack}(\upsilon, \chi)(T)$，则对每个
$t \in T$ 都有
$$
\chi(Y_t, (\text{id} \times f)^*\mathcal{M}^{\otimes n}) =
\chi(Y_t,
\mathcal{N}_t^{\otimes n} \otimes_{\mathcal{O}_{Y_t}}
f_t^*\mathcal{L}_t^{\otimes n}) =
\chi(Y_t, \mathcal{N}_t^{\otimes 2n}) = P(2n)
$$
其中在中间的等式中使用了同构
$\varphi : f^*\mathcal{L}_T \to \mathcal{N}_T$。
令 $P'(t) = P(2t)$，可知上述态射
$$
\mathit{Isom}_{\Polarizedstack}(\upsilon, \chi)
\longrightarrow
\mathit{Isom}_U(Y, X)
$$
的像包含在交
$$
\mathit{Isom}_U(Y, X) \cap \mathit{Mor}^{P', \mathcal{M}}_U(Y, X)
$$
中。该交是 $\mathit{Mor}_U(Y, X)$ 的两个开子空间之交
（见引理 \ref{moduli-lemma-Isom-in-Mor} 和注
\ref{moduli-remark-Mor-numerical}）。现在，
$\mathit{Mor}^{P', \mathcal{M}}_U(Y, X)$ 是 Noether 代数空间，
因为由引理 \ref{moduli-lemma-Mor-qc-over-base}，它在 $U$ 上有限表示。
所以该交也是 Noether 代数空间。由于态射
$$
\mathit{Isom}_{\Polarizedstack}(\upsilon, \chi)
\longrightarrow
\mathit{Isom}_U(Y, X) \cap \mathit{Mor}^{P', \mathcal{M}}_U(Y, X)
$$
仿射（见上文），结论成立。
\end{proof}

\begin{lemma}
\label{moduli-lemma-polarized-qs-lfp}
态射 $\Polarizedstack \to \Spec(\mathbf{Z})$ 拟分离且局部有限表示。
\end{lemma}

\begin{proof}
要检验 $\Polarizedstack \to \Spec(\mathbf{Z})$ 拟分离，需要证明其
对角态射拟紧且拟分离。这立即由引理
\ref{moduli-lemma-polarized-diagonal-separated-fp} 得出。
要证明 $\Polarizedstack \to \Spec(\mathbf{Z})$ 局部有限表示，
只需证明 $\Polarizedstack$ 保持极限，见《叠的极限》中的命题
\ref{stacks-limits-proposition-characterize-locally-finite-presentation}。
这就是《Quot》中的引理 \ref{quot-lemma-polarized-limits}。
\end{proof}

\begin{lemma}
\label{moduli-lemma-bounded-polarized}
设 $n \geq 1$ 是整数，$P$ 是数值多项式。设
$$
T \subset |\Polarizedstack|
$$
是具有下述性质的子集：对每个 $\xi \in T$，存在一个域 $k$ 和
$\Polarizedstack$ 在 $k$ 上的一个表示 $\xi$ 的对象
$(X, \mathcal{L})$，使得
\begin{enumerate}
\item $\mathcal{L}$ 在 $X$ 上的 Hilbert 多项式为 $P$；并且
\item 存在闭浸入 $i : X \to \mathbf{P}^n_k$，使得
$i^*\mathcal{O}_{\mathbf{P}^n}(1) \cong \mathcal{L}$。
\end{enumerate}
则 $T$ 是 Noether 拓扑空间，特别地，它拟紧。
\end{lemma}

\begin{proof}
注意，$|\Polarizedstack|$ 是局部 Noether 拓扑空间，见《叠的态射》
中的引理 \ref{stacks-morphisms-lemma-Noetherian-topology}
（这里还用到 $\Spec(\mathbf{Z})$ 是 Noether 的，因而由引理
\ref{moduli-lemma-polarized-qs-lfp} 和《叠的态射》中的引理
\ref{stacks-morphisms-lemma-locally-finite-type-locally-noetherian}，
$\Polarizedstack$ 是局部 Noether 代数叠）。因此，
$|\Polarizedstack|$ 的任意拟紧子集都是 Noether 拓扑空间，而这种
子集的任意子集也 Noether，见《拓扑》中的引理
\ref{topology-lemma-finite-union-Noetherian} 和
\ref{topology-lemma-Noetherian}。
% P12-ERR-0223: see ../control/ERRATA_CANDIDATES.jsonl
所以只需找到一个包含 $T$ 的拟紧子集。

\medskip\noindent
由引理 \ref{moduli-lemma-hilb-proper-over-base}，代数空间
$$
H =
\Hilbfunctor^{P, \mathcal{O}(1)}_{\mathbf{P}^n_\mathbf{Z}/\Spec(\mathbf{Z})}
$$
在 $\Spec(\mathbf{Z})$ 上真。由《Quot》中的引理
\ref{quot-lemma-extend-hilb-to-spaces}\footnote{
% P12-ERR-0224: see ../control/ERRATA_CANDIDATES.jsonl
我们将在后文看到（在此插入未来引用），$H$ 是概形，因而无需使用
本引理和《Quot》中的引理
\ref{quot-lemma-extend-polarized-to-spaces}。}，
$H$ 的恒等态射对应于闭子空间
$$
Z \subset \mathbf{P}^n_H
$$
它在 $H$ 上真、平坦且有限表示；限制
$\mathcal{N} = \mathcal{O}(1)|_Z$ 在 $Z/H$ 上相对丰富，并且在
$Z \to H$ 的纤维上具有 Hilbert 多项式 $P$。特别地，偶
$(Z \to H, \mathcal{N})$ 定义态射
$$
H \longrightarrow \Polarizedstack
$$
它把概形态射 $U \to H$ 送到族
$(Z_U \to U, \mathcal{N}_U)$ 的分类态射，见《Quot》中的引理
\ref{quot-lemma-extend-polarized-to-spaces}。
% P12-ERR-0225: see ../control/ERRATA_CANDIDATES.jsonl
由于 $H$ 是 Noether 代数空间（因为它在 $\mathbf{Z})$ 上真），
$|H|$ 是 Noether 的，因而拟紧。映射
$$
|H| \longrightarrow |\Polarizedstack|
$$
连续，所以其像拟紧。因此只需证明 $T$ 包含于
$|H| \to |\Polarizedstack|$ 的像中。然而，假设 (1) 和 (2)
恰好表达了这一事实：
% P12-ERR-0226: see ../control/ERRATA_CANDIDATES.jsonl
任取闭浸入 $i : X \to \mathbf{P}^n_k$，满足
$i^*\mathcal{O}_{\mathbf{P}^n}(1) \cong \mathcal{L}$，
由 $H$ 的模解释可得到 $H$ 的一个 $k$-值点。
这就完成了本引理的证明。
\end{proof}





\section{真态射上复形模空间的性质}
\label{moduli-section-complexes}

\noindent
设 $f : X \to B$ 是真、平坦且有限表示的代数空间态射。则参数化
负自 Ext 消失的相对完美复形的叠 $\Complexesstack_{X/B}$ 是代数的。
见《Quot》中的定理 \ref{quot-theorem-complexes-algebraic}。

\begin{lemma}
\label{moduli-lemma-complexes-diagonal-affine-fp}
$\Complexesstack_{X/B}$ 在 $B$ 上的对角态射仿射且有限表示。
\end{lemma}

\begin{proof}
《Quot》中的引理 \ref{quot-lemma-complexes-diagonal} 已证明该对角态射
可由代数空间表示。从其证明可知，需要证明下述事实：给定 $B$ 上概形
$T$ 以及对象 $E, E' \in D(\mathcal{O}_{X_T})$，使得
$(T, E)$ 和 $(T, E')$ 是 $\Complexesstack_{X/B}$ 在 $T$ 上的
纤维范畴中的对象，则
$\mathit{Isom}(E, E') \to T$ 仿射且有限表示。这里
$\mathit{Isom}(E, E')$ 是函子
$$
(\Sch/T)^{opp} \to \textit{Sets},\quad
T' \mapsto \{\varphi : E_{T'} \to E'_{T'}
\text{ 是下列范畴中的同构：}D(\mathcal{O}_{X_{T'}})\}
$$
其中 $E_{T'}$ 和 $E'_{T'}$ 是 $E$ 和 $E'$ 向 $X_{T'}$ 的导出拉回。
考虑由规则
$$
(\Sch/T)^{opp} \to \textit{Sets},\quad
% P12-ERR-0227: see ../control/ERRATA_CANDIDATES.jsonl
T' \mapsto \Hom_{\mathcal{O}_{X_{T'}}}(E_T, E'_T)
$$
定义的函子 $H = \SheafHom(E, E')$。由《Quot》中的引理
\ref{quot-lemma-complexes-open-neg-exts-vanishing}，这是在 $T$
上仿射且有限表示的代数空间。对
$H' = \SheafHom(E', E)$、$I = \SheafHom(E, E)$ 和
$I' = \SheafHom(E', E')$ 也同样。因此可见
$$
% P12-ERR-0228: see ../control/ERRATA_CANDIDATES.jsonl
\mathit{Isom}(E, E') = (H' \times_T H) \times_{c, I \times_T I', \sigma} T
$$
其中
$c(\varphi', \varphi) = (\varphi \circ \varphi', \varphi' \circ \varphi)$，
而 $\sigma = (\text{id}, \text{id})$
（与《Quot》中的命题 \ref{quot-proposition-isom} 的证明比较）。
因此，作为在 $T$ 上仿射的概形的纤维积，
$\mathit{Isom}(E, E')$ 在 $T$ 上仿射。类似地，
$\mathit{Isom}(E, E')$ 在 $T$ 上有限表示。
\end{proof}

\begin{lemma}
\label{moduli-lemma-complexes-qs-lfp}
态射 $\Complexesstack_{X/B} \to B$ 拟分离且局部有限表示。
\end{lemma}

\begin{proof}
要检验 $\Complexesstack_{X/B} \to B$ 拟分离，需要证明其对角态射
拟紧且拟分离。这立即由引理
\ref{moduli-lemma-complexes-diagonal-affine-fp} 得出。
要证明 $\Complexesstack_{X/B} \to B$ 局部有限表示，需要证明
$\Complexesstack_{X/B} \to B$ 保持极限，见《叠的极限》中的命题
\ref{stacks-limits-proposition-characterize-locally-finite-presentation}。
这由《Quot》中的引理 \ref{quot-lemma-complexes-limits} 得出
（略去一个小细节）。
\end{proof}
